\documentclass[11pt,reqno]{amsart}

\usepackage[margin=1.1in]{geometry}
\usepackage{amsmath,amssymb,amsthm,amsfonts}
\usepackage{mathrsfs}
\usepackage{microtype}
\usepackage{enumitem}
\usepackage{xcolor}
\usepackage[colorlinks=true,linkcolor=blue!40!black,citecolor=blue!40!black,urlcolor=blue!40!black]{hyperref}

% Theorem environments
\theoremstyle{plain}
\newtheorem{theorem}{Theorem}[section]
\newtheorem{proposition}[theorem]{Proposition}
\newtheorem{lemma}[theorem]{Lemma}
\newtheorem{corollary}[theorem]{Corollary}

\theoremstyle{definition}
\newtheorem{definition}[theorem]{Definition}
\newtheorem{remark}[theorem]{Remark}

% Notation
\newcommand{\NN}{\mathbb{N}}
\newcommand{\ZZ}{\mathbb{Z}}
\newcommand{\RR}{\mathbb{R}}
\newcommand{\PP}{\mathbb{P}}
\newcommand{\EE}{\mathbb{E}}
\newcommand{\Var}{\operatorname{Var}}
\DeclareMathOperator{\li}{Li}
\newcommand{\logstar}{\log_{\!*}}
\newcommand{\eps}{\varepsilon}

\title{Normal orders of the Erd\H{o}s chain functions $h(n)$ and $H(n)$}

\author{David Turturean\\
  Massachusetts Institute of Technology\\
  \texttt{davidct@mit.edu}}
\date{April 26, 2026}

\begin{document}

\begin{abstract}
For a positive integer $n$ let $h(n)$ be the largest $\ell$ such that there
exist primes $p_1<\cdots<p_\ell$ all dividing $n$ with $p_{i+1}\equiv 1\pmod{p_i}$,
and let $H(n)$ be the analogous quantity in which the $p_i$ are replaced by
arbitrary divisors $d_i\mid n$. Erd\H{o}s asked
\cite{ErdosProblems696} for estimates of $h(n)$ and $H(n)$, and whether
$H(n)/h(n)\to\infty$ for almost all $n$; he believed that the normal order of
$h(n)$ should be $\logstar n$.

We prove that for all but $o(x)$ integers $n\le x$,
\begin{equation*}
h(n) = \Bigl(\tfrac12 + o(1)\Bigr)\logstar n,
\qquad
H(n) = \bigl(1+o(1)\bigr)\logstar n,
\qquad
\frac{H(n)}{h(n)} = 2 + o(1).
\end{equation*}
In particular $H(n)/h(n)\to\infty$ is \emph{false} on a density-one set, while
Erd\H{o}s's belief about the order of magnitude of $h(n)$ is correct but with the
explicit normal-order constant $1/2$.

The upper bounds use a bad-pair density argument combined with an iterated-tower
counting lemma. The lower bound for $h$ is a Siegel--Walfisz prime-successor
construction transferred to integers via the Chinese remainder theorem. The
lower bound for $H$ uses a new \emph{subset-product successor lemma}: if
$g_1,\ldots,g_K$ are independent uniformly distributed elements of a finite
abelian group $G$ of order $N$, then with probability at least
$1-N/(2^K-1)$ some nonempty subproduct of the $g_i$ is equal to the identity.
This is applied, after a near-uniform residue reduction using Siegel--Walfisz,
inside Mertens-sized prime windows.
\end{abstract}

\maketitle

\section{Introduction}

Let $n\ge 2$ be an integer. Erd\H{o}s introduced two natural ``chain'' functions
attached to the multiplicative structure of $n$:
\begin{itemize}[leftmargin=2em]
\item[(i)] $h(n)$ is the largest integer $\ell$ such that there exist primes
$p_1<p_2<\cdots<p_\ell$ all dividing $n$, with
$p_{i+1}\equiv 1 \pmod{p_i}$ for $1\le i<\ell$;
\item[(ii)] $H(n)$ is the largest integer $u$ such that there exist integers
$d_1<d_2<\cdots<d_u$ all dividing $n$, with
$d_{i+1}\equiv 1 \pmod{d_i}$ for $1\le i<u$.
\end{itemize}
We set $h(1)=H(1)=0$. For $n\ge 2$, $h(n)\ge 1$ because $n$ has a prime
divisor, and $H(n)\ge 1$ because the one-term divisor chain $1\mid n$ is
allowed.

Problem 696 of the Erd\H{o}s Problems database \cite{ErdosProblems696} asks: estimate
$h(n)$ and $H(n)$, and decide whether
\begin{equation}\label{eq:erdos-conjecture}
\frac{H(n)}{h(n)} \;\longrightarrow\; \infty
\qquad\text{for almost all }n.
\end{equation}
According to the Erd\H{o}s Problems Database \cite{ErdosProblems696}, Erd\H{o}s
further suggested that the \emph{normal order} of $h(n)$ should be the iterated
logarithm $\logstar n$. (We were unable to locate this conjecture in
Erd\H{o}s's published papers; the reader should treat the attribution as based
on the database commentary.)

The first nontrivial result on this problem (folklore; see van Doorn's comments
in \cite{ErdosProblems696}) is that $h(n)\to\infty$ for almost all $n$. The
present paper resolves the full problem by determining the normal orders of
both $h$ and $H$ \emph{with explicit constants}.

\subsection{Statement of the result}\label{ss:main-result}

Define $T_0:=e$ and $T_{m+1}:=\exp(T_m)$ (these are formally introduced in
Section~\ref{sec:notation}). We set
\begin{equation}\label{eq:logstar-def}
\logstar x \;:=\; \min\{m\ge 0 : x\le T_m\}.
\end{equation}
Equivalently, for $x\ge e$, this is
\[
\logstar x \;=\; \min\{m\ge 0 : \log^{(m)} x \le e\}.
\]
Thus $\logstar e=0$, $\logstar(e^e)=1$, etc. This convention differs from
some other conventions by at most $1$, which is immaterial for the
normal-order statements below.

\begin{theorem}[Main theorem]\label{thm:main}
For all but $o(x)$ integers $n\le x$,
\begin{align}
h(n) &= \Bigl(\tfrac12 + o(1)\Bigr)\logstar n, \label{eq:h-normal}\\
H(n) &= \bigl(1 + o(1)\bigr)\logstar n, \label{eq:H-normal}\\
\frac{H(n)}{h(n)} &= 2 + o(1). \label{eq:ratio}
\end{align}
In particular, the divergence \eqref{eq:erdos-conjecture} is false on a
density-one set: $H(n)/h(n)$ is bounded above by $2+o(1)$ for almost all $n$.
\end{theorem}

It is striking that the constant $1/2$ for $h$ and the constant $1$ for $H$ both
appear, and that the ratio is exactly $2$. The factor of $2$ has a transparent
combinatorial explanation: a divisor-chain step $d_i\to d_{i+1}$ with
$d_{i+1}\equiv 1\pmod{d_i}$ allows $d_{i+1}$ to be any product of selected primes
hitting the trivial residue class mod $d_i$, whereas a prime-chain step $p_i\to
p_{i+1}$ requires $p_{i+1}$ itself to be a prime in the residue class
$1\pmod{p_i}$. For divisor chains the cost of one tower step is essentially
$\log_2 d_i$ random primes; for prime chains the cost is essentially $d_i$ itself,
and the bad-pair argument shows that two prime steps lower the iterated-tower
rank by exactly two when one would have been lowered by one for a single
divisor step.

\subsection{Comparison with what was previously known}

The qualitative claim $h(n)\to\infty$ a.a.\ was known. The quantitative
statement $h(n)=\Theta(\logstar n)$ a.a.\ does not appear to have been published
formally, although the proof we give in Section~\ref{sec:upper-h} for the upper
bound is essentially classical, drawing on Brun--Titchmarsh and Mertens, and the
$\Theta$-form lower bound follows the strategy implicit in
Erd\H{o}s--Granville--Pomerance--Spiro \cite{EGPS1990} on iterates of the
totient function.

The closest published work in print is Ford, Konyagin and Luca \cite{FKL2010},
who study the Pratt-tree depth
\[
  H_{\mathrm{Pr}}(p) \;:=\; \max\bigl\{k : \exists\,p_1\prec p_2\prec\cdots\prec p_k = p\bigr\},
\]
where each $p_i$ is prime and $a\prec b$ means $b\equiv 1\pmod a$. They prove
$H_{\mathrm{Pr}}(p)\le (\log p)^{1-c'}$ for almost all primes $p\le x$ (some
$c'>0$), and conjecture that the normal order is $e\log\log x$ via a branching
random walk model. \emph{This is a different problem.} In our setup, the chain
primes/divisors must all divide a specified integer $n$, which is a much more
restrictive constraint than ``terminating at a fixed prime $p$''. The Pratt-tree
depth grows like $\log\log p$; our $h(n)$ grows like $\frac12\logstar n$, much
slower, because the supply of chain primes is restricted to the
$\sim\log\log n$ prime divisors of $n$.

The sharp normal orders \eqref{eq:h-normal}--\eqref{eq:H-normal} (with the
matching constants $1/2$ and $1$) and the resulting refutation
\eqref{eq:ratio} of \eqref{eq:erdos-conjecture}, are, to the best of the authors'
knowledge, new. The upper bound for $h$ in
\eqref{eq:h-normal} is a Brun--Titchmarsh refinement of the bad-pair argument; the
matching lower bound is built greedily on the
Siegel--Walfisz prime-counting function in arithmetic progressions. The
lower bound \eqref{eq:H-normal} for $H$ rests on a probabilistic
\emph{subset-product successor lemma}, stated and proved in
Section~\ref{sec:subset-product}; this is the new technical input of the paper.

\subsection{Acknowledgement of AI assistance}

The arguments in this paper were assembled in collaboration with the Anthropic
language model Claude Opus, working from a four-turn proof sketch produced by
GPT-5.5-Pro (extended) on April 26, 2026. The authors verified all
displayed inequalities and, in several places, used GPT-5.5-Pro (extended
thinking) to fill in arguments that the original sketch had recorded only
schematically. The result statement and overall proof structure are due to
GPT-5.5-Pro; the iterative patches refining the writeup are also due to
GPT-5.5-Pro (extended), assembled together with Claude Code.

\subsection{Organization}

Section~\ref{sec:notation} fixes notation and recalls the analytic inputs.
Section~\ref{sec:upper-H} proves the upper bound $H(n)\le(1+o(1))\logstar n$ a.a.
Section~\ref{sec:upper-h} proves the upper bound $h(n)\le(\tfrac12+o(1))\logstar n$ a.a.
Section~\ref{sec:lower-h} proves the matching lower bound $h(n)\ge(\tfrac12-o(1))\logstar n$ a.a.\ via Siegel--Walfisz prime mass.
Section~\ref{sec:subset-product} establishes the subset-product successor lemma.
Section~\ref{sec:lower-H} proves $H(n)\ge(1-o(1))\logstar n$ a.a.\ using that lemma.
Section~\ref{sec:proof-of-main} assembles the main theorem.
Section~\ref{sec:remarks} contains concluding remarks.

\section{Notation, conventions, and analytic inputs}\label{sec:notation}

\subsection{Asymptotic conventions}

Throughout the paper $x$ is a real parameter tending to $\infty$. We write $f
\ll g$ or $f=O(g)$ to mean that $|f|\le C|g|$ for some absolute constant $C>0$
and all sufficiently large $x$, and $f=o(g)$ if $f/g\to 0$ as $x\to\infty$. The
implicit constants in any $O$ or $\ll$ may depend only on absolute parameters
introduced \emph{before} the relevant inequality is asserted, and are explicitly
labeled where needed.

We say a property $\mathcal{P}=\mathcal{P}(n)$ holds \emph{for almost all $n$}
(equivalently, ``a.a.''), or that $\mathcal{P}$ holds with \emph{density one},
if
\begin{equation}\label{eq:density-one}
\#\bigl\{n\le x : \neg\mathcal{P}(n)\bigr\} \;=\; o(x)
\qquad (x\to\infty).
\end{equation}

\subsection{Iterated logarithms and towers}

All logarithms in the paper are natural logarithms. We use $\log^{(m)}x$ for
the $m$-fold iterated logarithm, with $\log^{(0)}x=x$.

We use the iterated logarithm $\logstar x$ as in \eqref{eq:logstar-def}.
Equivalently, the tower of base-$e$ exponentials is defined by
\begin{equation}\label{eq:tower}
T_0 \;=\; e, \qquad T_{m+1} \;=\; \exp(T_m) \qquad (m\ge 0),
\end{equation}
and $\logstar x\le m$ if and only if $x\le T_m$. Note the basic identity
\begin{equation}\label{eq:logT}
\log T_m \;=\; T_{m-1} \qquad (m\ge 1).
\end{equation}

We will use the auxiliary scale
\begin{equation}\label{eq:Um}
U_m \;:=\; T_m^3 \qquad (m\ge 0),
\end{equation}
which behaves comparably to the tower $T_m$ but provides slack for
constant-factor losses.

We record several elementary computations that are used repeatedly:
\begin{align}
\log U_m &= 3 T_{m-1}, \label{eq:logU1}\\
\log\log U_m &= \log\bigl(3T_{m-1}\bigr) = T_{m-2} + \log 3, \label{eq:logU2}\\
\log\log\log U_m &= \log\bigl(T_{m-2}+\log 3\bigr) = T_{m-3} + O\bigl(T_{m-3}^{-1}\bigr).\label{eq:logU3}
\end{align}
Identities \eqref{eq:logU1}--\eqref{eq:logU3} are valid for $m\ge 4$, by repeated use of \eqref{eq:logT}.

\subsection{Standard analytic inputs}

We will appeal to four classical theorems of analytic number theory; we quote them in the precise form used.

\begin{lemma}[Siegel--Walfisz, {\cite[Ch.~22]{Davenport}, \cite{SiegelWalfisz}}]\label{lem:SW}
For every $A_0>0$ there exists $c_0=c_0(A_0)>0$ such that, uniformly for
$t\ge 3$, $q\le(\log t)^{A_0}$, and $(a,q)=1$,
\begin{equation}\label{eq:SW}
\pi(t;q,a)
\;=\;
\frac{\li(t)}{\varphi(q)}
+
O_{A_0}\!\left(t\exp\!\bigl(-c_0\sqrt{\log t}\bigr)\right).
\end{equation}
Consequently, if $3\le X\le Y$, $q\le(\log X)^{A_0}$, and $(a,q)=1$, then
\begin{equation}\label{eq:SW-reciprocal}
\sum_{\substack{X<p\le Y\\ p\equiv a\pmod q}}\frac1p
\;=\;
\frac1{\varphi(q)}\int_X^Y\frac{dt}{t\log t}
+
O_{A_0}\!\left(\exp\!\bigl(-c_0\sqrt{\log X}\bigr)\right),
\end{equation}
after possibly decreasing $c_0$. The implied constants depend only on $A_0$.
\end{lemma}

\begin{proof}[Proof of \eqref{eq:SW-reciprocal} given \eqref{eq:SW}]
Since $q\le(\log X)^{A_0}$ and $t\ge X$, \eqref{eq:SW} is valid uniformly
throughout $[X,Y]$. Partial summation gives
\[
\sum_{\substack{X<p\le Y\\ p\equiv a\pmod q}}\frac1p
\;=\;
\frac{\pi(Y;q,a)}{Y}
-\frac{\pi(X;q,a)}{X}
+
\int_X^Y \frac{\pi(t;q,a)}{t^2}\,dt.
\]
Substituting \eqref{eq:SW}, the main term is
\[
\frac1{\varphi(q)}
\left(
\frac{\li(Y)}{Y}
-\frac{\li(X)}{X}
+\int_X^Y\frac{\li(t)}{t^2}\,dt
\right)
\;=\;
\frac1{\varphi(q)}\int_X^Y\frac{dt}{t\log t}
\]
(after partial integration of the right-hand side, using $\li'(t)=1/\log t$).
The error term is
\[
O_{A_0}\!\left(
\exp(-c_0\sqrt{\log X})
+
\int_X^Y \frac{\exp(-c_0\sqrt{\log t})}{t}\,dt
\right).
\]
With $u=\sqrt{\log t}$, the integral is
\[
\ll \int_{\sqrt{\log X}}^\infty u\,e^{-c_0u}\,du
\;\ll\; \exp\!\left(-\tfrac{c_0}{2}\sqrt{\log X}\right),
\]
so, after decreasing $c_0$, the total error is
$O_{A_0}(\exp(-c_0\sqrt{\log X}))$. \qed
\end{proof}

\begin{remark}[Uniformity in the Siegel--Walfisz input]\label{rem:SW-uniformity}
The constant $c_0$ and all implied constants in Lemma~\ref{lem:SW}
depend only on the fixed exponent $A_0$. In the sequel we use this lemma
only with $A_0=1$; after this choice, the constants are absolute. In particular,
when \eqref{eq:SW-reciprocal} is applied on an interval $(X,Y]$, the error term
is uniform in the upper endpoint $Y$, because the proof bounded the error
integral by the tail integral over $[X,\infty)$. The usual ineffectivity of the
Siegel--Walfisz constant is irrelevant here, since all constants are fixed
before $x\to\infty$ and the arguments are qualitative density-one estimates.
\end{remark}

\begin{remark}[Amount of Siegel--Walfisz actually needed]\label{rem:SW-strength-needed}
The exponential error term in Lemma~\ref{lem:SW} is convenient, but the
arguments below do not rely on optimizing this error. It would be enough, in
all later applications, to have the following weaker consequence for one fixed
large exponent $B_{\mathrm{SW}}$:
\begin{equation}\label{eq:SW-weak-suffices}
\sum_{\substack{X<p\le Y\\ p\equiv a\pmod q}}\frac1p
=
\frac1{\varphi(q)}\int_X^Y\frac{dt}{t\log t}
+
O_{A_0,B_{\mathrm{SW}}}\bigl((\log X)^{-B_{\mathrm{SW}}}\bigr),
\end{equation}
uniformly for $3\le X\le Y$, $q\le(\log X)^{A_0}$, and $(a,q)=1$.
Indeed, in Lemma~\ref{lem:prime-successor-mass} one has
$X=\exp(\exp y)$, so the error in \eqref{eq:SW-weak-suffices} would be
$O(\exp(-B_{\mathrm{SW}}y))$, which is negligible compared with the main term
there. In Lemma~\ref{lem:block-eq} one has $X=\alpha\ge \exp y$ and
$d\le y$; after summing over at most $\varphi(d)\le y$ reduced residue classes,
the contribution of the error term would be
$O(y^{1-B_{\mathrm{SW}}})$, which is $O(y^{-1/2})$ as soon as
$B_{\mathrm{SW}}\ge 3$. Thus the later proof uses only standard
Siegel--Walfisz uniformity for fixed powers of $\log X$; the stronger
exponential form is quoted to keep the notation clean.
\end{remark}

\begin{lemma}[Brun--Titchmarsh, {\cite[Thm.~6.6]{IwaniecKowalski}, \cite{MontgomeryVaughan}}]\label{lem:BT}
There is an absolute constant $C_{BT}>0$ such that for all $q\ge 1$, all $a$ coprime to $q$, and all $t\ge 2q$,
\begin{equation}\label{eq:BT}
\pi(t;q,a) \;\le\; \frac{C_{BT}\, t}{\varphi(q)\,\log(t/q)}.
\end{equation}
\end{lemma}

\begin{lemma}[Mertens, {\cite[Thm.~427]{HardyWright}, \cite{Mertens-orig}}]\label{lem:Mertens}
There is an absolute constant $\mathcal{M}\in\RR$ such that
\begin{equation}\label{eq:Mertens}
\sum_{p\le t} \frac{1}{p} \;=\; \log\log t \;+\; \mathcal{M} \;+\; O\!\left(\frac{1}{\log t}\right)
\qquad (t\ge 2).
\end{equation}
\end{lemma}

\begin{lemma}[Chebyshev, {\cite[Thm.~414]{HardyWright}}]\label{lem:Cheb}
There is an absolute constant $C_{\theta}>0$ such that
\begin{equation}\label{eq:Cheb}
\sum_{p\le t} \log p \;\le\; C_{\theta}\, t \qquad (t\ge 2).
\end{equation}
\end{lemma}

\subsection{Probabilistic primes}\label{ss:prob-model}

Throughout the lower-bound arguments we work in the
\emph{independent prime-divisibility model}: we declare each prime $p$ to be
``selected'' independently with probability $1/p$. We then transfer back to
genuine integers $n\le x$ via the Chinese remainder theorem.

The transfer principle is the following standard observation, recorded once for
later reuse:

\begin{lemma}[CRT transfer]\label{lem:CRT}
Let $P\ge 2$, $M=\prod_{p\le P}p$, and let $\mathcal{E}$ be an event depending
only on the indicator vector $\bigl(\mathbf{1}_{p\mid n}\bigr)_{p\le P}$. Let
$\PP_{\mathrm{prod}}(\mathcal{E})$ denote the probability of $\mathcal{E}$ in
the product model where the $\mathbf{1}_{p\mid n}$ are independent
Bernoulli$(1/p)$ random variables. Then
\begin{equation}\label{eq:CRT}
\frac{1}{x}\#\{n\le x : \mathcal{E}(n) \text{ holds}\} \;=\; \PP_{\mathrm{prod}}(\mathcal{E})\;+\;O\!\bigl(M/x\bigr).
\end{equation}
In particular, if $M=o(x)$, the density of integers $n\le x$ for which
$\mathcal{E}$ holds equals $\PP_{\mathrm{prod}}(\mathcal{E})+o(1)$.
\end{lemma}

\begin{proof}[Sketch]
By the Chinese remainder theorem the map $n\bmod M\mapsto (n\bmod p)_{p\le P}$
is a bijection. Under the uniform distribution on $\ZZ/M\ZZ$, the events
$\{p\mid n\}$ for $p\le P$ are therefore independent and each has probability
$1/p$, exactly matching the product model. Hence the event $\mathcal{E}$ is a
union of residue classes mod $M$, with measure
$\PP_{\mathrm{prod}}(\mathcal{E})$ under the uniform distribution on $\ZZ/M\ZZ$.
Counting integers in $[1,x]$ in each residue class incurs an error of $\pm 1$
per class; since there are at most $M$ classes the total error is $O(M)$, and
dividing by $x$ gives \eqref{eq:CRT}.
\end{proof}

When the event $\mathcal{E}$ depends not only on prime divisibilities but also
on whether $n$ is divisible by some \emph{higher} prime power $p^2$, the same
argument applies with $M$ replaced by $\prod_{p\le P}p^{k_p}$. We will not need
this generalization.

\begin{remark}[Squarefree products under the CRT transfer]\label{rem:CRT-products}
In the lower-bound applications of Lemma~\ref{lem:CRT}, the constructed chain
terms are squarefree products of primes whose selected status is among the
coordinates $\bigl(\mathbf 1_{p\mid n}\bigr)_{p\le P}$. After transfer to an
actual integer $n$, ``selected'' means $p\mid n$. Hence any squarefree product
$e$ of selected primes with all prime factors at most $P$ satisfies $e\mid n$.
In particular, whenever such an $e$ occurs for an integer $n\le x$, one
automatically has $e\le n\le x$. Moreover $e\le M=\prod_{p\le P}p$, and in both
lower-bound applications below we verify $M=o(x)$; thus all squarefree products
considered in the CRT transfer are $<x$ for all sufficiently large $x$. No
additional size condition is hidden in the passage from the product model to
integers.
\end{remark}

\begin{lemma}[Chernoff bound for independent Bernoulli variables]\label{lem:chernoff}
Let $I_1,\ldots,I_M$ be independent $\{0,1\}$-valued random variables, not
necessarily identically distributed, and put
\[
W:=\sum_{k=1}^M I_k,\qquad \mu:=\EE W.
\]
Then, for every $0<\delta<1$,
\begin{equation}\label{eq:chernoff}
\PP\bigl(W\le (1-\delta)\mu\bigr)
\le
\exp\!\left(-\frac{\delta^2\mu}{2}\right).
\end{equation}
In particular,
\begin{equation}\label{eq:chernoff-one-third}
\PP\bigl(W\le (2/3)\mu\bigr)\le \exp(-\mu/18).
\end{equation}
\end{lemma}

\begin{proof}
Write $p_k:=\EE I_k$. For $\lambda>0$, Markov's inequality and independence give
\[
\PP(W\le (1-\delta)\mu)
=
\PP\bigl(e^{-\lambda W}\ge e^{-\lambda(1-\delta)\mu}\bigr)
\le
e^{\lambda(1-\delta)\mu}
\prod_{k=1}^M \bigl(1-p_k+p_k e^{-\lambda}\bigr).
\]
Using $1+u\le e^u$,
\[
\prod_{k=1}^M \bigl(1-p_k+p_k e^{-\lambda}\bigr)
\le
\exp\!\left(-(1-e^{-\lambda})\sum_{k=1}^M p_k\right)
=
\exp\!\left(-(1-e^{-\lambda})\mu\right).
\]
Choose $\lambda=\log(1/(1-\delta))$, so that $1-e^{-\lambda}=\delta$. Then
\[
\PP(W\le (1-\delta)\mu)
\le
\exp\!\left(
-\mu\left[\delta-(1-\delta)\log\frac1{1-\delta}\right]
\right).
\]
The elementary inequality
\[
\delta-(1-\delta)\log\frac1{1-\delta}\ge \frac{\delta^2}{2}
\qquad (0<\delta<1)
\]
gives \eqref{eq:chernoff}. Taking $\delta=1/3$ gives
\eqref{eq:chernoff-one-third}.
\end{proof}

\begin{remark}[CRT transfer for $x$-dependent events]\label{rem:CRT-x-dependent}
Lemma~\ref{lem:CRT} is applied below to events $\mathcal E=\mathcal E_x$ and
cutoffs $P=P(x)$ that vary with $x$. This causes no additional difficulty:
for each fixed $x$, the lemma is an exact finite CRT statement with error
$O(M/x)$, where $M=\prod_{p\le P}p$. Thus it is enough, in each lower-bound
application, to verify $M=o(x)$ for the particular cutoff used there.

In practice we verify the stronger condition $P=o(\log x)$. By
Lemma~\ref{lem:Cheb},
\[
\log M=\sum_{p\le P}\log p\le C_\theta P=o(\log x),
\]
and therefore $M=\exp(o(\log x))=x^{o(1)}=o(x)$. The deterministic rules used
later, such as choosing the least selected prime or the least admissible
squarefree product, ensure that the relevant events are genuine functions of
the finite indicator vector
\[
\bigl(\mathbf 1_{p\mid n}\bigr)_{p\le P}.
\]
No information about primes above $P$ or about higher prime powers is used in
those CRT transfers.
\end{remark}

\section{Upper bound for $H(n)$}\label{sec:upper-H}

In this section we prove:

\begin{theorem}\label{thm:upper-H}
For all but $o(x)$ integers $n\le x$,
\begin{equation}\label{eq:thm-upper-H}
H(n) \;\le\; \bigl(1 + o(1)\bigr)\logstar x.
\end{equation}
\end{theorem}

The strategy: any divisor chain step $d_i<d_{i+1}$ with $d_{i+1}\equiv 1\pmod{d_i}$ that does not satisfy $d_{i+1}\ge \exp(\sqrt{d_i})$ creates a ``bad pair''
$(d_i,d_{i+1})$. We show that bad pairs of moderate size are rare on average,
so for almost all $n$ no large step is bad, forcing each chain step to climb
the iterated tower by roughly one full level.

\subsection{The bad-pair lemma}

\begin{definition}[Bad pair, divisor version]\label{def:bad-pair-H}
An ordered pair of positive integers $(d,e)$ is called
\emph{H-bad} if
\begin{equation}\label{eq:bad-pair-def-H}
d<e, \qquad e\equiv 1 \pmod{d}, \qquad e<\exp\bigl(\sqrt{d}\bigr).
\end{equation}
\end{definition}

\begin{lemma}\label{lem:bad-pair-density-H}
Let $Y=Y(x)\to\infty$. The number of integers $n\le x$ such that some
$H$-bad pair $(d,e)$ with $d\ge Y$ has both $d\mid n$ and $e\mid n$ is at most
\begin{equation}\label{eq:bad-pair-count-H}
O\!\left(x\, Y^{-1/2}\right) \;=\; o(x).
\end{equation}
\end{lemma}

\begin{proof}
The congruence $e\equiv 1\pmod d$ forces $\gcd(d,e)=1$, hence
$\mathrm{lcm}(d,e)=de$ and the number of $n\le x$ divisible by both
$d$ and $e$ is $\lfloor x/(de)\rfloor\le x/(de)$. The total count is therefore
bounded by
\begin{equation}\label{eq:bad-pair-sum1-H}
\sum_{d\ge Y} \frac{x}{d}
\sum_{\substack{d<e<\exp(\sqrt d)\\ e\equiv 1\pmod d}} \frac{1}{e}.
\end{equation}
For fixed $d\ge 4$, write $e=ad+1$ with $1\le a\le \exp(\sqrt d)/d$:
\begin{equation}\label{eq:inner-sum-H}
\sum_{\substack{d<e<\exp(\sqrt d)\\ e\equiv 1\pmod d}}\frac{1}{e}
\;\le\;\sum_{1\le a\le \exp(\sqrt d)/d} \frac{1}{ad+1}
\;\le\; \frac{1}{d}\sum_{1\le a\le \exp(\sqrt d)/d} \frac{1}{a}.
\end{equation}
The harmonic sum is $\le 1+\log\bigl(\exp(\sqrt d)/d\bigr)\le \sqrt d$, valid for
$d\ge 4$. Substituting,
\begin{equation}\label{eq:inner-sum-H2}
\sum_{\substack{d<e<\exp(\sqrt d)\\ e\equiv 1\pmod d}}\frac{1}{e}
\;\le\; \frac{\sqrt d}{d}
\;=\; d^{-1/2}.
\end{equation}
Inserting \eqref{eq:inner-sum-H2} into \eqref{eq:bad-pair-sum1-H},
\begin{equation}\label{eq:bad-pair-sum2-H}
\sum_{d\ge Y}\frac{x}{d}\cdot d^{-1/2}
\;=\; x\sum_{d\ge Y} d^{-3/2}
\;\ll\; x\, Y^{-1/2}.
\end{equation}
Choosing $Y\to\infty$ gives \eqref{eq:bad-pair-count-H}.
\end{proof}

\subsection{The U-tower lemma for $H$}

Set
\begin{equation}\label{eq:G-def}
G(t) \;:=\; (\log t)^2.
\end{equation}

\begin{lemma}\label{lem:U-tower-H}
For every $m\ge 5$,
\begin{equation}\label{eq:U-tower-H}
G(U_m) \;\le\; U_{m-1}.
\end{equation}
\end{lemma}

\begin{proof}
By \eqref{eq:logU1}, $\log U_m = 3T_{m-1}$, hence
\[
G(U_m) \;=\; (\log U_m)^2 \;=\; 9 T_{m-1}^2.
\]
We must show $9T_{m-1}^2\le T_{m-1}^3=U_{m-1}$, equivalently $T_{m-1}\ge 9$. Since
$T_4=\exp\bigl(\exp\bigl(\exp(e)\bigr)\bigr)$ is enormous, this holds whenever
$m-1\ge 1$, i.e.\ $m\ge 2$; we asked $m\ge 5$ to absorb the $\log 3$ in
\eqref{eq:logU2} for any later use, but it is not needed for
\eqref{eq:U-tower-H}.
\end{proof}

\subsection{Proof of Theorem~\ref{thm:upper-H}}

Let $x$ be large, $L:=\logstar x$, and put $Y:=\lfloor L^{1/2}\rfloor$.
By Lemma~\ref{lem:bad-pair-density-H}, all but $o(x)$ integers $n\le x$ have the
property
\begin{equation}\label{eq:no-H-bad}
\text{no $H$-bad pair $(d,e)$ with $d\ge Y$ has both $d\mid n$ and $e\mid n$}.
\tag{$\star_H$}
\end{equation}
Fix such an $n$ and let $d_1<d_2<\cdots<d_u$ be a divisor chain witnessing
$H(n)=u$. Let
$u_{\mathrm{small}}:=\#\{i:d_i<Y\}$ and $u_{\mathrm{large}}:=u-u_{\mathrm{small}}$.

Trivially $u_{\mathrm{small}}\le Y$, since the small terms are distinct integers in $[1,Y)$.

For each large index $i$ with $d_i\ge Y$ and $i<u$, the pair
$(d_i,d_{i+1})$ satisfies $d_i<d_{i+1}$, $d_{i+1}\equiv 1\pmod{d_i}$, both
divide $n$, and $d_i\ge Y$. Property \eqref{eq:no-H-bad} rules out
$d_{i+1}<\exp(\sqrt{d_i})$, so
\begin{equation}\label{eq:large-step-H}
d_{i+1} \;\ge\; \exp\bigl(\sqrt{d_i}\bigr),
\qquad\text{equivalently,}\qquad
d_i \;\le\; (\log d_{i+1})^2 \;=\; G(d_{i+1}).
\end{equation}

Define the \emph{$U$-rank} of $d$ to be
$\rho(d):=\min\{m\ge 0 : d\le U_m\}$. Since $d_u\le n\le x\le T_L\le U_L$, we
have $\rho(d_u)\le L$.

\begin{lemma}[Tower descent for $H$]\label{lem:tower-descent-H}
If $d_i\ge Y$ and $i<u$ then $\rho(d_i)\le \rho(d_{i+1})-1$, provided
$\rho(d_{i+1})\ge 6$.
\end{lemma}

\begin{proof}
Set $m:=\rho(d_{i+1})$. Then $d_{i+1}\le U_m$, so by \eqref{eq:large-step-H} and
Lemma~\ref{lem:U-tower-H},
\begin{equation}\label{eq:tower-descent-H-comp}
d_i \;\le\; G(d_{i+1}) \;\le\; G(U_m) \;\le\; U_{m-1},
\end{equation}
provided $m\ge 5$ so that $U_m\ge T_5\ge 9$. Hence $\rho(d_i)\le m-1=\rho(d_{i+1})-1$.
\end{proof}

For completeness, we spell out the rank count used in the next paragraph.
The terms with $d_i\ge Y$ form a suffix of the chain, since the chain is
strictly increasing. Also $\rho(d_i)$ is nondecreasing in $i$. If there is no
index $i$ with $d_i\ge Y$ and $\rho(d_i)\ge 6$, there is nothing to count.
Otherwise let $i_0$ be the first such index. Then for every $i_0\le i<u$ we
have $d_i\ge Y$ and
\[
\rho(d_{i+1})\ge \rho(d_i)\ge 6,
\]
so Lemma~\ref{lem:tower-descent-H} gives
\[
\rho(d_{i+1})\ge \rho(d_i)+1.
\]
Therefore
\[
\rho(d_u)\ge \rho(d_{i_0})+(u-i_0)\ge 6+(u-i_0).
\]
Thus the number of large terms with rank at least $6$ is
\[
u-i_0+1\le \rho(d_u)-5\le L.
\]

The number of large terms with $\rho(d_i)\ge 6$ is therefore at most $\rho(d_u)\le L$. The remaining large terms have $\rho(d_i)\le 5$, hence $d_i\le U_5$, hence there are $O(1)$ of them.

Putting everything together,
\begin{equation}\label{eq:H-final-bound}
H(n) \;=\; u_{\mathrm{small}}+u_{\mathrm{large}}
\;\le\; Y \;+\; L \;+\; O(1)
\;=\; L \;+\; O(L^{1/2})
\;=\; (1+o(1))L.
\end{equation}
This is exactly \eqref{eq:thm-upper-H}. \qed

\section{Upper bound for $h(n)$}\label{sec:upper-h}

In this section we prove:

\begin{theorem}\label{thm:upper-h}
For all but $o(x)$ integers $n\le x$,
\begin{equation}\label{eq:thm-upper-h}
h(n) \;\le\; \Bigl(\tfrac12 + o(1)\Bigr)\logstar x.
\end{equation}
\end{theorem}

The qualitative strategy is the same as for $H$, but the analogue of
the bad-pair lemma is sharper because $p,q$ must both be primes. The cost of
``$q\equiv 1\pmod p$ and $q$ prime'' is roughly $1/(\varphi(p)\log q)$ rather
than $1/q$, and a Brun--Titchmarsh estimate makes the bad-prime-pair tail
small enough to push the cutoff $\exp(\sqrt p)$ all the way up to
$Q(p):=\exp\bigl(\exp(p/(\log p)^2)\bigr)$. This faster threshold means each
prime-chain step descends \emph{two} levels of the iterated tower, halving
the count.

\subsection{The bad-prime-pair lemma}

For a prime $p\ge 3$ define
\begin{equation}\label{eq:Q-def}
Q(p) \;:=\; \exp\!\Bigl(\exp\!\Bigl(\frac{p}{(\log p)^2}\Bigr)\Bigr).
\end{equation}

\begin{definition}[Bad prime pair]\label{def:bad-pair-h}
An ordered pair of primes $(p,q)$ is called \emph{$h$-bad} if
\begin{equation}\label{eq:bad-pair-def-h}
p<q, \qquad q\equiv 1 \pmod p, \qquad q<Q(p).
\end{equation}
\end{definition}

\begin{lemma}\label{lem:bad-pair-density-h}
Let $Y\to\infty$. The number of $n\le x$ such that some $h$-bad pair $(p,q)$
with $p\ge Y$ has both $p\mid n$ and $q\mid n$ is $o(x)$.
\end{lemma}

\begin{proof}
For fixed prime $p$ we estimate
\begin{equation}\label{eq:bad-prime-pair-S}
\mathcal{S}(p) \;:=\; \sum_{\substack{p<q\le Q(p)\\ q\equiv 1\pmod p\\ q\text{ prime}}}\frac{1}{q}.
\end{equation}
Split this into the ranges $p<q\le p^2$ and $p^2<q\le Q(p)$.

\emph{Range $p<q\le p^2$.} Here $q$ is one of at most $p$ primes (in fact, the
trivial bound $q=p+1,\,2p+1,\dots$ gives at most $p-1$ candidates), so
\begin{equation}\label{eq:bad-h-low}
\sum_{\substack{p<q\le p^2\\ q\equiv 1\pmod p\\ q\text{ prime}}}\frac{1}{q}
\;\le\; \sum_{a=1}^{p-1}\frac{1}{ap+1}
\;\le\; \frac{1}{p}\sum_{a=1}^{p-1}\frac1a
\;\ll\; \frac{\log p}{p}.
\end{equation}

\emph{Range $p^2<q\le Q(p)$.} By Brun--Titchmarsh (Lemma~\ref{lem:BT}), for
$t\ge 2p$,
\[
\pi(t;p,1) \;\le\; \frac{C_{BT}\,t}{\varphi(p)\log(t/p)}
\;\le\; \frac{2C_{BT}\,t}{p\log(t/p)},
\]
since $\varphi(p)=p-1\ge p/2$ for $p\ge 2$. Partial summation yields
\begin{align}
\sum_{\substack{p^2<q\le Q(p)\\ q\equiv 1\pmod p\\ q\text{ prime}}}\frac{1}{q}
&= \int_{p^2}^{Q(p)} \frac{d\pi(t;p,1)}{t}\nonumber\\
&\le \frac{\pi(Q(p);p,1)}{Q(p)} + \int_{p^2}^{Q(p)}\frac{\pi(t;p,1)}{t^2}\,dt\nonumber\\
&\ll\; \frac{1}{p}\int_{p^2}^{Q(p)} \frac{dt}{t\log(t/p)}\nonumber\\
&=\;\frac{1}{p}\Bigl[\log\log(t/p)\Bigr]_{p^2}^{Q(p)}\nonumber\\
&\ll\; \frac{1}{p}\Bigl(\log\log Q(p) - \log\log p\Bigr).\label{eq:BT-pre}
\end{align}
In the Stieltjes partial summation formula the lower boundary contribution
$-\pi(p^2;p,1)/p^2$ is nonpositive and has therefore been discarded. The upper
boundary term is harmless: by Brun--Titchmarsh it is
\[
\frac{\pi(Q(p);p,1)}{Q(p)}
\ll
\frac{1}{p\log(Q(p)/p)},
\]
which is absorbed by the displayed integral bound.
By definition \eqref{eq:Q-def}, $\log\log Q(p)= p/(\log p)^2$, so
\begin{equation}\label{eq:BT-final}
\sum_{\substack{p^2<q\le Q(p)\\ q\equiv 1\pmod p\\ q\text{ prime}}}\frac{1}{q}
\;\ll\; \frac{1}{p}\cdot\frac{p}{(\log p)^2}
\;=\; \frac{1}{(\log p)^2}.
\end{equation}
Combining \eqref{eq:bad-h-low} and \eqref{eq:BT-final},
\begin{equation}\label{eq:Sp-bound}
\mathcal{S}(p) \;\ll\; \frac{\log p}{p} + \frac{1}{(\log p)^2}.
\end{equation}

For each $h$-bad pair $(p,q)$ with $p\ge Y$, the count of $n\le x$ with
$pq\mid n$ is $\le x/(pq)$ (recalling $\gcd(p,q)=1$). Summing,
\begin{align}
\#\{n\le x\colon\exists\,(p,q)\text{ $h$-bad}, p\ge Y, pq\mid n\}
&\le\; x\sum_{p\ge Y}\frac{\mathcal{S}(p)}{p}\nonumber\\
&\ll\; x\sum_{p\ge Y}\Bigl(\frac{\log p}{p^2}+\frac{1}{p(\log p)^2}\Bigr).\label{eq:bad-pair-sum-h}
\end{align}
Both series converge: $\sum_p (\log p)/p^2<\infty$ and $\sum_p 1/(p(\log p)^2)<\infty$.
Indeed, for $p\ge 3$ both prime sums are dominated by the corresponding
integer sums:
\[
\sum_{p\ge 3}\frac{\log p}{p^2}
\le
\sum_{n\ge 3}\frac{\log n}{n^2}
<\infty,
\qquad
\sum_{p\ge 3}\frac{1}{p(\log p)^2}
\le
\sum_{n\ge 3}\frac{1}{n(\log n)^2}
<\infty.
\]
The second integer sum converges by the integral test.
Therefore the right-hand side of \eqref{eq:bad-pair-sum-h} is $o(x)$
as $Y\to\infty$.
\end{proof}

\subsection{The U-tower lemma for $h$}

Define
\begin{equation}\label{eq:f-def}
f(t) \;:=\; C_f\,(\log\log t)\bigl(\log\log\log t\bigr)^2,
\end{equation}
where $C_f$ is a fixed absolute constant chosen large enough below.

\begin{lemma}\label{lem:U-tower-h}
There exists an integer $m_0^*=m_0^*(C_f)\ge 6$ such that, for all $m\ge m_0^*$,
\begin{equation}\label{eq:U-tower-h}
f(U_m) \;\le\; U_{m-2}.
\end{equation}
\end{lemma}

\begin{proof}
By \eqref{eq:logU2}, $\log\log U_m = T_{m-2}+\log 3$, and by \eqref{eq:logU3},
$\log\log\log U_m = T_{m-3}+O(T_{m-3}^{-1})$. Therefore, for $m\ge 5$,
\begin{align}
f(U_m) &= C_f\,(T_{m-2}+\log 3)\bigl(T_{m-3}+O(T_{m-3}^{-1})\bigr)^2\nonumber\\
&\le 2C_f\, T_{m-2}\,T_{m-3}^2,\label{eq:fUm}
\end{align}
provided $T_{m-3}\ge 100$ (take $m\ge 6$). Now $T_{m-3}=\log T_{m-2}$, so
$T_{m-3}^2=(\log T_{m-2})^2$. We must compare this with $T_{m-2}^2$:
\[
\frac{f(U_m)}{U_{m-2}}\;\le\;\frac{2C_f\,T_{m-2}\,T_{m-3}^2}{T_{m-2}^3}\;=\;\frac{2C_f\,(\log T_{m-2})^2}{T_{m-2}^2}.
\]
Since $(\log s)^2/s^2\to 0$ as $s\to\infty$, this ratio tends to $0$, so
$f(U_m)/U_{m-2}\le 1$ for all sufficiently large $m$. Concretely, choose
$s_0=s_0(C_f)$ so large that
\[
   \frac{2C_f(\log s)^2}{s^2}\le 1 \qquad (s\ge s_0),
\]
and then choose $m_0^*$ so that $m_0^*\ge 6$ and
$T_{m_0^*-2}\ge s_0$. This yields \eqref{eq:U-tower-h}.
\end{proof}

\begin{remark}[Monotonicity of the inversion function]\label{rem:f-monotone}
We will use that the function
\[
f(t)=C_f(\log\log t)(\log\log\log t)^2
\]
is increasing for all sufficiently large $t$, in particular for $t\ge T_3$.
Indeed, writing $s=\log\log t$, we have
\[
\frac{1}{C_f}f(t)=s(\log s)^2.
\]
For $s>1$,
\[
\frac{d}{ds}\bigl(s(\log s)^2\bigr)
=
(\log s)^2+2\log s>0,
\]
and $s=\log\log t$ is itself increasing in $t$. Since
$\log\log T_3=T_1>1$, the asserted monotonicity holds for $t\ge T_3$.
\end{remark}

\subsection{Proof of Theorem~\ref{thm:upper-h}}

Let $x$ be large, $L:=\logstar x$, $Y:=\lfloor L^{1/2}\rfloor$. By
Lemma~\ref{lem:bad-pair-density-h}, all but $o(x)$ integers $n\le x$ satisfy
\begin{equation}\label{eq:no-h-bad}
\text{no $h$-bad pair $(p,q)$ with $p\ge Y$ has both $p\mid n$ and $q\mid n$}.
\tag{$\star_h$}
\end{equation}
Fix such an $n$, and let $p_1<p_2<\cdots<p_\ell$ be a prime chain witnessing
$h(n)=\ell$. Let $\ell_{\mathrm{small}}=\#\{i:p_i<Y\}\le Y$, and $\ell_{\mathrm{large}}=\ell-\ell_{\mathrm{small}}$.

For each large index $i$ with $p_i\ge Y$ and $i<\ell$, the pair $(p_i,p_{i+1})$ is
not $h$-bad, so
\begin{equation}\label{eq:large-step-h}
p_{i+1} \;\ge\; Q(p_i)
\;=\; \exp\!\Bigl(\exp\!\Bigl(\tfrac{p_i}{(\log p_i)^2}\Bigr)\Bigr).
\end{equation}
Solving \eqref{eq:large-step-h} for $p_i$:
\[
\frac{p_i}{(\log p_i)^2} \;\le\; \log\log p_{i+1},
\quad\text{so}\quad
p_i \;\le\; (\log p_i)^2\,\log\log p_{i+1}.
\]
We extract a clean algebraic inversion. Set $a:=\log p_i$, $b:=\log\log p_{i+1}$;
the inequality above reads $e^a\le a^2 b$.

Since $p_i\ge Y$ and $Y\to\infty$, \eqref{eq:large-step-h} gives
\[
   b=\log\log p_{i+1}\ge \frac{p_i}{(\log p_i)^2}\longrightarrow\infty,
\]
uniformly for the large-index steps under consideration. Thus all estimates
below that require $b$ sufficiently large apply after increasing $x$ if
necessary.

More explicitly, after increasing $x$ if necessary we have $Y>e^2$, and the
function $t\mapsto t/(\log t)^2$ is increasing for $t>e^2$. Hence throughout
the large-index part of the chain,
\[
   b=\log\log p_{i+1}
   \ge \frac{p_i}{(\log p_i)^2}
   \ge \frac{Y}{(\log Y)^2}
   \longrightarrow \infty .
\]
Thus the threshold beyond which the elementary inversion below is valid may be
chosen independently of the particular chain and of the index $i$.

We claim
\begin{equation}\label{eq:a-bound}
a \;\le\; \log b\;+\;3\log\log b
\end{equation}
for all sufficiently large $b$. Indeed, the function $u\mapsto e^u/u^2$ is
increasing for $u>2$. If $a>\log b+3\log\log b$, then
\[
\frac{e^a}{a^2}
\;>\;
\frac{\exp(\log b+3\log\log b)}{(\log b+3\log\log b)^2}
\;=\;
b\cdot\frac{(\log b)^3}{(\log b+3\log\log b)^2}
\;>\;b
\]
for all sufficiently large $b$, contradicting $e^a\le a^2 b$. This establishes
\eqref{eq:a-bound}. Therefore $\log p_i = a\le \log b + 3\log\log b$, and so
\begin{equation}\label{eq:p_i-bound}
p_i \;\le\; (\log p_i)^2\cdot b
\;=\; a^2\cdot b
\;\le\; (\log b + 3\log\log b)^2\cdot b
\;\le\; C_f\,b\,(\log b)^2,
\end{equation}
for an absolute constant $C_f>0$ (take $C_f=16$, say) and all sufficiently
large $b$. Recalling $b=\log\log p_{i+1}$, we obtain
\begin{equation}\label{eq:p_i-final}
p_i \;\le\; C_f\,(\log\log p_{i+1})\,\bigl(\log\log\log p_{i+1}\bigr)^2 \;=\; f(p_{i+1}).
\end{equation}

We now run the same $U$-rank descent as in Section~\ref{sec:upper-H}. Let
$\rho(p):=\min\{m\ge 0:p\le U_m\}$.

\begin{lemma}[Tower descent for $h$]\label{lem:tower-descent-h}
Provided $\rho(p_{i+1})\ge m_0^*$ (where $m_0^*$ is from Lemma~\ref{lem:U-tower-h})
and $p_i\ge Y\ge 100$, $i<\ell$,
\begin{equation}\label{eq:tower-descent-h}
\rho(p_i) \;\le\; \rho(p_{i+1})\;-\;2.
\end{equation}
\end{lemma}

\begin{proof}
Set $m:=\rho(p_{i+1})\ge m_0^*$. Then $p_{i+1}\le U_m$. The function
$f(t)=C_f(\log\log t)(\log\log\log t)^2$ is monotone increasing for
$t\ge T_3$, hence by \eqref{eq:p_i-final},
\[
p_i \;\le\; f(p_{i+1})\;\le\; f(U_m)\;\le\; U_{m-2},
\]
the last inequality being Lemma~\ref{lem:U-tower-h}. Hence $\rho(p_i)\le m-2$.
\end{proof}

Again we record the counting argument explicitly. The terms with $p_i\ge Y$
form a suffix of the prime chain, and $\rho(p_i)$ is nondecreasing in $i$.
If there is no index $i$ with $p_i\ge Y$ and $\rho(p_i)\ge m_0^*$, then there
are no large high-rank terms to count. Otherwise let $i_0$ be the first such
index. For every $i_0\le i<\ell$ we have $p_i\ge Y$ and
\[
\rho(p_{i+1})\ge \rho(p_i)\ge m_0^*,
\]
so Lemma~\ref{lem:tower-descent-h} gives
\[
\rho(p_{i+1})\ge \rho(p_i)+2.
\]
Consequently
\[
\rho(p_\ell)\ge \rho(p_{i_0})+2(\ell-i_0)\ge m_0^*+2(\ell-i_0).
\]
Thus the number of large terms with rank at least $m_0^*$ is at most
\[
\ell-i_0+1
\le
1+\frac{\rho(p_\ell)-m_0^*}{2}
\le
\frac{L}{2}+O(1).
\]

Since $p_\ell\le n\le x\le T_L\le U_L$, we have $\rho(p_\ell)\le L$. The number
of large terms $i$ with $\rho(p_i)\ge m_0^*$ is therefore at most $L/2+O(1)$,
since each such large term is at least 2 levels below the next. The remaining
large terms have $\rho(p_i)<m_0^*$, hence $p_i<U_{m_0^*-1}=O(1)$, so there are
$O(1)$ of them.

Adding $\ell_{\mathrm{small}}\le Y=\lfloor L^{1/2}\rfloor$,
\begin{equation}\label{eq:h-final-bound}
h(n) \;=\; \ell \;\le\; \tfrac{L}{2} + O(1) + Y
\;=\; \tfrac{L}{2} + O(L^{1/2})
\;=\; \bigl(\tfrac12 + o(1)\bigr) L,
\end{equation}
proving \eqref{eq:thm-upper-h}. \qed

\section{Lower bound for $h(n)$}\label{sec:lower-h}

We now prove the matching lower bound:

\begin{theorem}\label{thm:lower-h}
For all but $o(x)$ integers $n\le x$,
\begin{equation}\label{eq:thm-lower-h}
h(n) \;\ge\; \Bigl(\tfrac12 - o(1)\Bigr)\logstar x.
\end{equation}
\end{theorem}

The strategy is to construct, for ``random'' $n$, a long prime chain $p_1<\cdots<p_R$ greedily: at each stage we use Siegel--Walfisz to ensure that the next prime in the residue class $1\bmod p_j$, lying inside a controlled tower window, exists with overwhelming probability. We then transfer to integers via the CRT.

Throughout this section we put
\[
   L:=\logstar x.
\]
All choices of scales below are functions of this parameter $L$, and hence of
$x$.

\subsection{The prime-successor mass lemma}

\begin{lemma}[Uniform prime-successor mass]\label{lem:prime-successor-mass}
There are absolute constants $c_0>0$, $y_0\ge 1$ such that the following
holds. For every $y\ge y_0$, every prime $p\le y$, and every $\eta\ge 2$,
\[
\sum_{\substack{\exp(\exp y)<q\le \exp(\exp(\eta y))\\
q\equiv 1\pmod p\\ q\ {\rm prime}}}\frac1q
\;=\;
\frac{(\eta-1)y}{\varphi(p)}
\;+\;
O\!\left(\exp(-c_0\exp(y/2))\right),
\]
uniformly in $p,y,\eta$. In particular,
\[
\sum_{\substack{\exp(\exp y)<q\le \exp(\exp(\eta y))\\
q\equiv 1\pmod p\\ q\ {\rm prime}}}\frac1q
\;\ge\;
\frac12\,\frac{(\eta-1)y}{p}
\]
for all sufficiently large $y$.
\end{lemma}

\begin{proof}
Let
\[
Z=\exp(\exp y),\qquad X=\exp(\exp(\eta y)).
\]
For $t\ge Z$, $\log t\ge \exp y$, hence $p\le y\le(\log Z)$ for all
sufficiently large $y$, so Lemma~\ref{lem:SW} with $A_0=1$ applies on $[Z,X]$.
By \eqref{eq:SW-reciprocal},
\[
\sum_{\substack{Z<q\le X\\ q\equiv 1\pmod p}}\frac1q
\;=\;
\frac1{\varphi(p)}\int_Z^X\frac{dt}{t\log t}
\;+\;
O\!\left(\exp(-c_0\sqrt{\log Z})\right).
\]
The integral evaluates to
$\log\log X-\log\log Z=(\eta-1)y$, and $\sqrt{\log Z}=\exp(y/2)$, so the
error is $O(\exp(-c_0\exp(y/2)))$. This proves the asymptotic formula.

Since $p\le y$, $\varphi(p)=p-1\le p$, and $(\eta-1)y/p\ge 1$ for $\eta\ge 2$
and $p\le y$. For sufficiently large $y$, the error term is therefore smaller
than half the main lower bound, giving
\[
\sum_{\substack{Z<q\le X\\ q\equiv 1\pmod p}}\frac1q
\;\ge\;
\frac12\,\frac{(\eta-1)y}{p}. \qedhere
\]
\end{proof}

\subsection{Tower scale for $h$ and the greedy construction}

Fix $B:=L^2$ and the integers
\begin{equation}\label{eq:m0-mj-h}
m_0 \;=\; \lfloor L^{1/2}\rfloor,
\qquad
m_j \;=\; m_0+2(j-1) \quad (j\ge 1),
\end{equation}
so that $m_j$ increments by 2. Let $R$ be the largest integer with $m_R\le L-2$;
then
\begin{equation}\label{eq:R-h}
R \;=\; \frac{L-m_0-2}{2} + O(1) \;=\; \Bigl(\tfrac12 - o(1)\Bigr)L.
\end{equation}

Define
\begin{equation}\label{eq:y_j-h}
y_j \;=\; T_{m_j}/B \qquad (1\le j\le R+1).
\end{equation}

Since $T_{m_{j+1}}=T_{m_j+2}=\exp\bigl(\exp(T_{m_j})\bigr)=\exp(\exp(B y_j))$, we have
\begin{equation}\label{eq:y_j+1-vs-y_j}
y_{j+1} \;=\; \frac{T_{m_{j+1}}}{B}\;=\;\frac{\exp(\exp(B y_j))}{B}.
\end{equation}
In particular,
\begin{equation}\label{eq:loglog-y_{j+1}}
\log\log y_{j+1} \;=\; \log\bigl(\exp(B y_j)-\log B\bigr)\;=\; By_j + O(1).
\end{equation}

\begin{lemma}\label{lem:greedy-h}
Work in the independent prime model on primes $r\le y_R$. The probability that
the following event $\mathcal{C}_h$ \emph{fails} is $o(1)$:
\begin{quote}
$\mathcal{C}_h$:\ \ there exist primes $p_1<\cdots<p_R$ all selected, with
$p_j\le y_j$ and $p_{j+1}\equiv 1\pmod{p_j}$ for $1\le j<R$.
\end{quote}
\end{lemma}

\begin{proof}
We expose primes in stages and make deterministic choices throughout. If there
is at least one selected prime $p\le y_1$, define $p_1$ to be the least
selected prime $\le y_1$; otherwise the construction fails at the initial
stage.

The probability of initial failure is
\begin{equation}\label{eq:greedy-init}
\prod_{p\le y_1}\!\bigl(1-1/p\bigr)
\;\le\;
\exp\!\Bigl(-\!\!\sum_{p\le y_1}\!1/p\Bigr)
\;=\;
\exp(-\log\log y_1+O(1)).
\end{equation}
Since $y_1=T_{m_0}/B$ and $m_0\to\infty$, we have
$\log\log y_1=T_{m_0-2}+O(\log B)\to\infty$, so \eqref{eq:greedy-init} is
$o(1)$.

Now suppose that $p_1,\dots,p_j$ have already been defined, with
$p_i\le y_i$ and $p_{i+1}\equiv 1\pmod{p_i}$ for $1\le i<j$. Define the
candidate set
\[
\mathcal Q_j(p_j)
\;:=\;
\bigl\{
q\ \text{prime}:
q\equiv 1\!\pmod{p_j},\
\exp(\exp y_j)<q\le \exp(\exp((B/2)y_j))
\bigr\}.
\]
If at least one prime in $\mathcal Q_j(p_j)$ is selected, define $p_{j+1}$ to
be the least selected prime in $\mathcal Q_j(p_j)$; otherwise the construction
fails at stage $j$.

By induction, every prime whose selection has been inspected before stage $j$
is at most $y_j$: the initial step inspects only primes $\le y_1$, and the
search interval at stage $r<j$ is contained in $(\exp(\exp y_r),y_{r+1}]$,
with $y_{r+1}\le y_j$. On the other hand every prime in
$\mathcal Q_j(p_j)$ is $>\exp(\exp y_j)>y_j$. Thus the primes inspected at
stage $j$ are disjoint from all previously inspected primes. This justifies
conditioning on the entire past while retaining independent Bernoulli
selection for the primes in $\mathcal Q_j(p_j)$.

For $L$ large, this search interval is contained in
$(\exp(\exp y_j),\,y_{j+1}]$, since by \eqref{eq:y_j+1-vs-y_j}
\[
\exp(\exp((B/2)y_j))
\;\ll\;
\frac{\exp(\exp(B y_j))}{B}
\;=\;
y_{j+1}.
\]
Hence, on success at stage $j$, $p_{j+1}\le y_{j+1}$. Also
$p_{j+1}>\exp(\exp y_j)>y_j\ge p_j$, so the chain remains strictly increasing.

Apply Lemma~\ref{lem:prime-successor-mass} with $y=y_j$, $\eta=B/2$, and
modulus $p=p_j\le y_j$. Since $B=L^2$, we have $\eta\ge 2$ for large $L$, so
\begin{equation}\label{eq:greedy-mass}
\sum_{q\in\mathcal Q_j(p_j)}\frac1q
\;\ge\;
\frac12\,\frac{(B/2-1)y_j}{p_j}
\;\ge\;
\frac{B}{8},
\end{equation}
for $B\ge 4$, since $p_j\le y_j$.

The construction up to stage $j$ depends only on selected primes lying below
the present search interval, whereas every prime in $\mathcal Q_j(p_j)$ is
$>\exp(\exp y_j)>y_j\ge p_j$. Hence, conditional on the past, the selection
events for primes in $\mathcal Q_j(p_j)$ are still independent Bernoulli events
with probabilities $1/q$. Therefore
\[
\PP(\text{stage $j$ fails}\mid\text{past})
\;\le\;
\prod_{q\in\mathcal Q_j(p_j)}\!\bigl(1-1/q\bigr)
\;\le\;
\exp\!\Bigl(-\!\!\!\sum_{q\in\mathcal Q_j(p_j)}\!\!\!1/q\Bigr)
\;\le\;
\exp(-B/8)
\;=\;
\exp(-L^2/8).
\]
Union-bounding over $1\le j<R\le L$,
\begin{equation}\label{eq:greedy-union}
\PP(\mathcal{C}_h^c) \;\le\; o(1) + L\cdot\exp(-L^2/8) \;=\; o(1). \qedhere
\end{equation}
\end{proof}

\subsection{Proof of Theorem~\ref{thm:lower-h}}

We apply Lemma~\ref{lem:CRT} with the cutoff $P:=y_R$. The event
$\mathcal{C}_h$ depends only on which primes $p\le y_R$ divide $n$, so it is a
union of residue classes mod $M:=\prod_{p\le y_R}p$. We need $M=o(x)$.

Indeed, the initial choice of $p_1$ only inspects primes $p\le y_1\le y_R$.
At stage $j<R$, the search interval $\mathcal Q_j(p_j)$ is contained in
$(\exp(\exp y_j),y_{j+1}]$, as shown above, and hence contains only primes
$\le y_R$. The deterministic rule ``choose the least selected prime'' makes the
event $\mathcal C_h$ a function only of the indicator vector
$\bigl(\mathbf 1_{p\mid n}\bigr)_{p\le y_R}$. No higher prime powers are used.

Since $m_R\le L-2$, $y_R\le T_{L-2}/B$. On the other hand, by the definition of
$L=\logstar x$ we have
\[
   T_{L-1}<x\le T_L
\]
for all sufficiently large $L$. Taking one logarithm in the lower inequality
and using $\log T_{L-1}=T_{L-2}$ gives
\begin{equation}\label{eq:logx-vs-T}
\log x \;>\; T_{L-2}.
\end{equation}
Hence $y_R\le T_{L-2}/L^2 = o(\log x)$. By Chebyshev (Lemma~\ref{lem:Cheb}),
$\log M=\sum_{p\le y_R}\log p\le C_\theta y_R = o(\log x)$, so $M=x^{o(1)}=o(x)$.

We also record explicitly that the constructed prime divisibilities are
compatible with the restriction $n\le x$. If $\mathcal C_h$ is realized by an
integer $n$, then each constructed prime $p_j$ is one of the primes whose
indicator is $1$, and hence $p_j\mid n$. Since the primes $p_1,\dots,p_R$ are
distinct and all lie below $P=y_R$, their product is at most
\[
\prod_{p\le y_R}p=M=o(x).
\]
Thus the CRT transfer is prescribing a squarefree divisor of size $<x$ for all
sufficiently large $x$; equivalently, by Remark~\ref{rem:CRT-products}, no
extra size obstruction is present.

By Lemma~\ref{lem:CRT}, the density of $n\le x$ for which $\mathcal{C}_h$ holds
is $\PP_{\mathrm{prod}}(\mathcal{C}_h)+O(M/x)=1-o(1)$. Whenever $\mathcal{C}_h$
holds, the chain witnesses $h(n)\ge R=(\tfrac12-o(1))L$. This proves
\eqref{eq:thm-lower-h}. \qed

\section{The subset-product successor lemma}\label{sec:subset-product}

This section contains the key probabilistic ingredient for the lower bound on $H(n)$. We work in two steps. First, an abstract statement about random elements of finite abelian groups (Lemma~\ref{lem:subset-product}). Then a number-theoretic application that produces composite-divisor successors (Lemma~\ref{lem:composite-successor}).

\subsection{Abstract subset-product lemma}

Let $G$ be a finite abelian group, written multiplicatively, with identity
$1_G$ and order $N=|G|$.

In the next lemma no generation hypothesis is imposed on the sampled elements:
they are only required to be independent and uniform. In the number-theoretic
application the residue classes coming from good blocks need not generate $G$.

\begin{lemma}\label{lem:subset-product}
Let $g_1,\dots,g_K$ be independent uniformly distributed elements of $G$. Define
\begin{equation}\label{eq:Z-def}
Z \;:=\; \#\Bigl\{\emptyset\ne S\subseteq\{1,\dots,K\}:\;\prod_{i\in S} g_i = 1_G\Bigr\}.
\end{equation}
Then
\begin{align}
\EE[Z] &\;=\; \frac{2^K-1}{N}, \label{eq:EZ}\\
\Var[Z] &\;\le\; \EE[Z], \label{eq:VarZ}\\
\PP(Z=0) &\;\le\; \frac{N}{2^K-1}.\label{eq:PZzero}
\end{align}
\end{lemma}

\begin{proof}
For each nonempty $S\subseteq[K]$, the random variable $X_S:=\prod_{i\in S}g_i$
is uniformly distributed on $G$ (by independence of the $g_i$ and translation
invariance), so $\PP(X_S=1_G)=1/N$. Summing,
\[
\EE[Z] \;=\; \sum_{\emptyset\ne S\subseteq[K]}\PP(X_S=1_G)
\;=\;\frac{2^K-1}{N},
\]
proving \eqref{eq:EZ}.

To bound the variance we show that for any two distinct nonempty subsets $S\ne T$
of $[K]$, the random variables $X_S$ and $X_T$ are \emph{independent}. Pick
\begin{equation}\label{eq:i0-def}
i_0\;\in\;S\,\triangle\,T,
\end{equation}
the symmetric difference (which is nonempty since $S\ne T$). Without loss of
generality $i_0\in S\setminus T$ (else swap $S,T$). Condition on the joint value
of $\{g_j: j\ne i_0\}$. Under this conditioning $X_T$ is determined (because
$i_0\notin T$), so $X_T$ takes a fixed value; meanwhile $X_S=g_{i_0}\cdot\prod_{i\in
S\setminus\{i_0\}}g_i$ depends on $g_{i_0}$, which under our conditioning is
uniformly distributed on $G$. Thus $X_S\mid \{g_j: j\ne i_0\}$ is uniform on $G$,
hence in particular conditionally independent of $X_T$. Marginalizing,
\begin{equation}\label{eq:XSXT-indep}
X_S \;\text{and}\; X_T \;\text{are independent, both uniform on $G$.}
\end{equation}
Equivalently, for any $x_0,y_0\in G$, conditioning on all $g_j$ with
$j\ne i_0$ gives
\[
\begin{aligned}
\PP(X_S=x_0,\ X_T=y_0)
&=
\EE\!\left[
   \mathbf 1_{\{X_T=y_0\}}\,
   \PP(X_S=x_0\mid \{g_j:j\ne i_0\})
 \right]  \\
&=
\frac1N\,\PP(X_T=y_0)
=
\frac1{N^2}.
\end{aligned}
\]
This gives the asserted independence for every distinct pair
$S\ne T$.
Therefore $\PP(X_S=X_T=1_G)=1/N^2$. Hence
\[
\EE[Z^2] \;=\; \sum_{S}\PP(X_S=1)+\sum_{S\ne T}\PP(X_S=X_T=1)
\;=\;\frac{2^K-1}{N}+\frac{(2^K-1)(2^K-2)}{N^2}.
\]
Subtracting $(\EE Z)^2=(2^K-1)^2/N^2$,
\begin{equation}\label{eq:Var-formula}
\Var Z \;=\; \frac{2^K-1}{N}-\frac{2^K-1}{N^2}
\;=\;\frac{(2^K-1)(N-1)}{N^2}
\;\le\; \frac{2^K-1}{N}\;=\;\EE Z,
\end{equation}
proving \eqref{eq:VarZ}.

Finally, by Chebyshev's inequality,
\[
\PP(Z=0) \;\le\; \PP\bigl(|Z-\EE Z|\ge \EE Z\bigr) \;\le\; \frac{\Var Z}{(\EE Z)^2}
\;\le\;\frac{1}{\EE Z} \;=\;\frac{N}{2^K-1},
\]
proving \eqref{eq:PZzero}.
\end{proof}

\begin{remark}[Near-uniform version]\label{rem:near-uniform}
For our application, the $g_i$ are not exactly uniform on $G$ but only close to
uniform in total variation. We record the precise reduction. Let $u_G$ denote
the uniform probability measure on $G$, and suppose that the random variables
$g_1,\dots,g_K$ are independent, with laws $\mu_i$ satisfying
\[
d_{\rm TV}(\mu_i,u_G)\le \eps \qquad (1\le i\le K).
\]
For each $i$, choose a maximal coupling of $g_i$ with a uniform random variable
$\widetilde g_i\sim u_G$ such that
\[
\PP(g_i\ne \widetilde g_i)=d_{\rm TV}(\mu_i,u_G)\le \eps,
\]
and choose these couplings independently for different $i$. Then
$\widetilde g_1,\dots,\widetilde g_K$ are independent and exactly uniform on
$G$, and by the union bound
\[
\PP\bigl((g_1,\dots,g_K)\ne(\widetilde g_1,\dots,\widetilde g_K)\bigr)
\le K\eps.
\]
Let $Z$ be the number of identity subset-products formed from the $g_i$, and
let $\widetilde Z$ be the corresponding number for the $\widetilde g_i$. The
events $\{Z=0\}$ and $\{\widetilde Z=0\}$ can differ only if at least one
$g_i$ differs from $\widetilde g_i$. Hence Lemma~\ref{lem:subset-product} gives
\begin{equation}\label{eq:PZzero-perturbed}
\PP(Z=0)
\le
\PP(\widetilde Z=0)+K\eps
\le
\frac{N}{2^K-1}+K\eps.
\end{equation}
More generally, if the $i$th total-variation gap is $\eps_i$, the additive
error is $\sum_{i=1}^K \eps_i$.
\end{remark}

\subsection{The composite-successor lemma}

In the next lemma the subscript $A$ in $O_A(\cdot)$ refers only to the fixed
absolute window parameter used later in Lemma~\ref{lem:composite-successor}.
In the application below we take $A=20$, so all such constants are absolute.

\begin{lemma}[Weighted residue equidistribution in a block]\label{lem:block-eq}
Let $d\le y$, let $G:=(\ZZ/d\ZZ)^\times$, and let
$\mathcal{B}=(\alpha,\beta]$ be one of the blocks constructed below in the
proof of Lemma~\ref{lem:composite-successor} (so $\alpha\ge\exp y$ and the
reciprocal-prime mass satisfies $\mu(\mathcal{B}):=\sum_{q\in\mathcal{B}}1/q
=1+O(1/y)$). For every reduced residue class $a\bmod d$,
\[
\sum_{\substack{q\in\mathcal{B}\\q\equiv a\pmod d}}\frac1q
\;=\;
\frac{\mu(\mathcal{B})}{\varphi(d)}
\;+\;
O_A\!\bigl(\exp(-c\sqrt y)\bigr),
\]
uniformly in $d\le y$, $a\in G$, and the block $\mathcal{B}$. The same
estimate holds with $1/q$ replaced by $1/(q-1)$. Consequently, conditional
on the event that exactly one prime in $\mathcal{B}$ is selected, the residue
class of that prime modulo $d$ has total variation distance
$O(y^{-1/2})$ from the uniform distribution on $G$.
\end{lemma}

\begin{proof}
The case $d=1$ is trivial, since $G$ has one element and the total-variation
distance to the uniform distribution is zero. We may assume $d\ge 2$. Since
every prime in $\mathcal B$ is $>\exp y>d$, each such prime is coprime to $d$
and therefore defines an element of $G$.

Since $d\le y\le\log\alpha$, Lemma~\ref{lem:SW} with $A_0=1$ applies
uniformly on $(\alpha,\beta]$. By \eqref{eq:SW-reciprocal},
\[
\sum_{\substack{\alpha<q\le\beta\\q\equiv a\pmod d}}\frac1q
\;=\;
\frac1{\varphi(d)}\int_\alpha^\beta\frac{dt}{t\log t}
\;+\;
O\!\bigl(\exp(-c_0\sqrt{\log\alpha})\bigr).
\]
Since $\log\alpha\ge y$, the error is $O(\exp(-c_0\sqrt y))$. Write
\[
I_{\mathcal B}:=\int_\alpha^\beta\frac{dt}{t\log t}.
\]
For each reduced residue class $a\bmod d$ we have
\[
\sum_{\substack{q\in\mathcal B\\q\equiv a\pmod d}}\frac1q
=
\frac{I_{\mathcal B}}{\varphi(d)}
+
O\!\bigl(\exp(-c_0\sqrt y)\bigr),
\]
uniformly in $d\le y$, $a$, and $\mathcal B$. Summing this estimate over all
$a\in(\ZZ/d\ZZ)^\times$ gives
\[
\mu(\mathcal B)
=
I_{\mathcal B}
+
O\!\bigl(\varphi(d)\exp(-c_0\sqrt y)\bigr).
\]
Dividing by $\varphi(d)$ and substituting back yields
\[
\frac{I_{\mathcal B}}{\varphi(d)}
=
\frac{\mu(\mathcal B)}{\varphi(d)}
+
O\!\bigl(\exp(-c_0\sqrt y)\bigr),
\]
and hence
\[
\sum_{\substack{q\in\mathcal B\\q\equiv a\pmod d}}\frac1q
=
\frac{\mu(\mathcal B)}{\varphi(d)}
+
O\!\bigl(\exp(-c_0\sqrt y)\bigr).
\]
After decreasing $c_0$ if necessary, this is the stated
$O_A(\exp(-c\sqrt y))$ estimate.

For the weight $1/(q-1)$, note that
\[
\sum_{q\in\mathcal{B}}\left|\frac1{q-1}-\frac1q\right|
\;\ll\;
\sum_{q>\exp y}\frac1{q^2}
\;\ll\;
\exp(-y),
\]
and the same bound holds after restricting to any residue class.

Now condition on the event $\mathcal G_{\mathcal B}$ that exactly one prime in
$\mathcal B$ is selected, and let $Q_{\mathcal B}$ denote that unique selected
prime. For $a\in G$, we have
\begin{align*}
\PP(Q_{\mathcal B}\equiv a\pmod d\mid \mathcal G_{\mathcal B})
&=
\frac{
\sum_{\substack{q\in\mathcal{B}\\q\equiv a\pmod d}}
\frac1q\prod_{r\in\mathcal{B},\ r\ne q}(1-1/r)
}{
\sum_{q\in\mathcal{B}}
\frac1q\prod_{r\in\mathcal{B},\ r\ne q}(1-1/r)
}  \\
&=
\frac{
\sum_{\substack{q\in\mathcal{B}\\q\equiv a\pmod d}}\frac1{q-1}
}{
\sum_{q\in\mathcal{B}}\frac1{q-1}
},
\end{align*}
because the common factor $\prod_{r\in\mathcal B}(1-1/r)$ cancels.

By the estimate just proved with the weight $1/(q-1)$, the numerator is
\[
\frac{\mu(\mathcal B)}{\varphi(d)}
+
O_A\bigl(\exp(-c\sqrt y)\bigr),
\]
while the denominator is
\[
\mu(\mathcal B)+O(\exp(-y))=1+O(1/y).
\]
Since $\mu(\mathcal B)=1+O(1/y)$, it follows that
\[
\PP(Q_{\mathcal B}\equiv a\pmod d\mid \mathcal G_{\mathcal B})
=
\frac1{\varphi(d)}
+
O\!\left(\frac{1}{y\varphi(d)}\right)
+
O_A\bigl(\exp(-c\sqrt y)\bigr).
\]
Consequently the total-variation distance from the uniform distribution on
$G$ is at most
\[
\frac12\sum_{a\in G}
\left|
\PP(Q_{\mathcal B}\equiv a\pmod d\mid \mathcal G_{\mathcal B})
-\frac1{\varphi(d)}
\right|
\ll
\frac1y+y\exp(-c\sqrt y)
\ll y^{-1/2},
\]
for all sufficiently large $y$. \qed
\end{proof}

\begin{remark}[Uniformity under conditioning on good blocks]\label{rem:block-uniformity}
The total-variation estimate in Lemma~\ref{lem:block-eq} is uniform in
$d\le y$, in the reduced residue class, and in the block $\mathcal B$ among
the blocks constructed in Lemma~\ref{lem:composite-successor}. Consequently,
if we condition on any fixed collection of block indices being good, the unique
selected primes in those blocks remain independent across the chosen blocks,
and each of their residue classes modulo $d$ is still $O(y^{-1/2})$-close to
uniform on $(\ZZ/d\ZZ)^\times$. This is because the blocks are disjoint, and
inside a chosen block the conditioning is only the event that exactly one prime
of that block is selected.

The same conclusion holds when the block indices are chosen as the first
$K$ good blocks. Indeed, conditioning on a realized value
$k_1<\cdots<k_K$ of these indices imposes only the events that
$\mathcal B_{k_i}$ is good and that certain disjoint intervening blocks are
not good. The latter events involve no primes from the blocks
$\mathcal B_{k_i}$, and within each $\mathcal B_{k_i}$ the only conditioning
is that exactly one prime is selected. Therefore the unique selected primes
from the realized good blocks remain independent, and their residue laws retain
the same uniform $O(y^{-1/2})$ total-variation bound.
\end{remark}

\begin{lemma}\label{lem:composite-successor}
There exist absolute constants $A>10$, $c>0$, and $y_0\ge 2$ with the following
property. Let $y\ge y_0$ and $1\le d\le y$. Work in the independent prime model
in which each prime $q$ is selected independently with probability $1/q$. Then,
with probability at least $1-O(y^{-c})$, uniformly in $d\le y$, there exists a
squarefree product $e$ of \emph{selected} primes lying entirely in
$\bigl(\exp y,\,\exp(y^{A-1})\bigr]$, such that
\begin{equation}\label{eq:composite-successor}
e\equiv 1\pmod d, \qquad e>d, \qquad e\le \exp(y^A).
\end{equation}
\end{lemma}

\begin{proof}
Set $A:=20$ (any sufficiently large constant works; the value $20$ comfortably
satisfies all numerical thresholds below), and take $y_0$ large enough for the
estimates below.

We first dispose of the case $d=1$. In this case the congruence condition
$e\equiv 1\pmod d$ is vacuous. By Mertens' theorem,
\[
\sum_{\exp y<q\le \exp(y^2)}\frac1q
=
\log\log(\exp(y^2))-\log\log(\exp y)+O(1/y)
=
\log y+O(1/y).
\]
Hence the probability that no prime in $(\exp y,\exp(y^2)]$ is selected is at
most
\[
\exp\!\left(-\sum_{\exp y<q\le \exp(y^2)}\frac1q\right)
=
\exp\bigl(-\log y+O(1/y)\bigr)
\ll y^{-1}.
\]
On the complementary event, any selected prime $q\in(\exp y,\exp(y^2)]$ gives
$e=q$, which is squarefree, satisfies $e>1=d$, satisfies $e\equiv 1\pmod 1$,
lies in $\bigl(\exp y,\exp(y^{A-1})\bigr]$, and obeys
$e\le \exp(y^2)\le \exp(y^A)$. Thus the lemma holds for $d=1$, after decreasing
$c$ if necessary.

We may therefore assume $d\ge 2$. Write $G:=(\ZZ/d\ZZ)^\times$ and
$N:=\varphi(d)\le d\le y$.

\medskip\noindent
\textbf{Step 1: Mertens reciprocal mass on the prime window.} By
Lemma~\ref{lem:Mertens},
\begin{equation}\label{eq:Mertens-window}
\sum_{\exp y<q\le \exp(y^{A-1})}\frac{1}{q}
\;=\; \log\log\bigl(\exp(y^{A-1})\bigr)-\log\log(\exp y)+O(1)
\;=\; (A-2)\log y + O(1).
\end{equation}

\medskip\noindent
\textbf{Step 2: Block decomposition.}
Recall $A=20$. Fix $\gamma=15$ and set
\begin{equation}\label{eq:M-blocks}
M \;:=\; \lfloor\gamma\log y\rfloor.
\end{equation}
Define real endpoints $\alpha_k$ by
\[
\log\log \alpha_k \;=\; \log y+k-1, \qquad 1\le k\le M+1.
\]
Thus $\alpha_1=\exp y$, and
\[
\alpha_{M+1}
\;=\;
\exp\!\left(y\,e^M\right)
\;\le\;
\exp\!\left(y^{\gamma+1}\right)
\;\le\;
\exp(y^{A-1})
\]
for all large $y$, since $\gamma+1=16<A-1=19$. Let
\[
\mathcal B_k \;:=\; \bigl(\alpha_k,\alpha_{k+1}\bigr]
\]
denote the set of primes in this interval. By Mertens' theorem
(Lemma~\ref{lem:Mertens}),
\begin{equation}\label{eq:block-mass}
\mu_k \;:=\; \sum_{q\in\mathcal B_k}\frac1q
\;=\;
\log\log\alpha_{k+1}-\log\log\alpha_k+O(1/\log\alpha_k)
\;=\;
1+O(1/y),
\end{equation}
uniformly in $k$. The total mass absorbed by the first $M$ blocks is
$M\cdot(1+O(1/y))=\gamma\log y\cdot(1+o(1))$, which fits inside the window
since $\gamma<A-2$.

\medskip\noindent
\textbf{Step 3: Good blocks.} Call a block $\mathcal{B}_k$ \emph{good} if exactly
one prime in $\mathcal{B}_k$ is selected. Put
\[
\mu_k:=\sum_{q\in\mathcal B_k}\frac1q=1+O(1/y),
\qquad
\Pi_k:=\prod_{q\in\mathcal B_k}\Bigl(1-\frac1q\Bigr).
\]
Then
\begin{align*}
\PP(\mathcal B_k\text{ is good})
&=
\sum_{q\in\mathcal B_k}
\frac1q\prod_{q'\in\mathcal B_k\setminus\{q\}}
\Bigl(1-\frac1{q'}\Bigr)  \\
&=
\Pi_k\sum_{q\in\mathcal B_k}\frac{1/q}{1-1/q}
=
\Pi_k\sum_{q\in\mathcal B_k}\frac1{q-1}.
\end{align*}
Moreover,
\[
\sum_{q\in\mathcal B_k}\frac1{q-1}
=
\sum_{q\in\mathcal B_k}\frac1q
+
O\!\left(\sum_{q>\exp y}\frac1{q^2}\right)
=
\mu_k+O(e^{-y})
=
1+O(1/y),
\]
and
\[
\log \Pi_k
=
\sum_{q\in\mathcal B_k}\log\Bigl(1-\frac1q\Bigr)
=
-\mu_k+O\!\left(\sum_{q>\exp y}\frac1{q^2}\right)
=
-1+O(1/y).
\]
Therefore
\[
\PP(\mathcal B_k\text{ is good})=e^{-1}+O(1/y).
\]
Hence, writing $\rho:=e^{-1}$,
\begin{equation}\label{eq:good-block-prob}
\PP(\mathcal{B}_k\text{ is good}) \;=\; \rho \;+\; O(1/y)
\qquad\text{(uniformly in $k$).}
\end{equation}

The events $\{\mathcal{B}_k\text{ is good}\}$ for different $k$ depend on
\emph{disjoint} sets of primes, so they are mutually independent.

\medskip\noindent
\textbf{Step 4: Chernoff for the number of good blocks.} Let
$I_k$ be the indicator of the event that $\mathcal B_k$ is good, and let
\[
W:=\sum_{k=1}^M I_k.
\]
The variables $I_k$ are independent, since the blocks use disjoint sets of
primes. By \eqref{eq:good-block-prob},
\[
\EE W
=
\sum_{k=1}^M \PP(\mathcal B_k\text{ is good})
=
\rho M+O(M/y)
=
\rho\gamma\log y+O(1).
\]
Choose
\begin{equation}\label{eq:gamma-choice}
\kappa \;:=\; \tfrac{\rho\gamma}{2}>\frac{1}{\log 2}+1,
\qquad
\text{i.e.\ } \gamma > 2(\rho^{-1})\!\bigl(\tfrac{1}{\log 2}+1\bigr).
\end{equation}
Numerically $\rho=e^{-1}\approx 0.368$, $1/\log 2\approx 1.443$, so it suffices
to take $\gamma\ge 14$. We also require $\gamma<A-2=18$, so
$\gamma\in[14,18)$ is admissible; we keep the choice $\gamma=15$.

For all sufficiently large $y$, the threshold $\kappa\log y$ is at most
$(2/3)\EE W$. The Chernoff bound for sums of independent Bernoulli variables
therefore gives
\begin{equation}\label{eq:lots-good-blocks}
\PP\bigl(W\;<\;\kappa\log y\bigr)
\le
\PP\bigl(W<(2/3)\EE W\bigr)
\le
\exp(-\EE W/18)
=
O(y^{-c_1})
\end{equation}
for some constant $c_1=c_1(\gamma)>0$.
The use of Chernoff here is justified by Lemma~\ref{lem:chernoff}: the
indicators $I_k$ are independent because the prime blocks are disjoint, and
they need not be identically distributed. The estimate uses only their
independence and the value of $\EE W$.

\medskip\noindent
\textbf{Step 5: Conditional residue distribution.} Put
\[
K:=\lceil \kappa\log y\rceil.
\]
On the event $W\ge K$, let
\[
k_1<k_2<\cdots<k_K
\]
be the first $K$ indices for which $\mathcal B_{k_i}$ is good. From each good
block $\mathcal B_{k_i}$ we extract the unique selected prime $q_{k_i}$, which
lies in $\mathcal B_{k_i}\subset(\exp y,\exp(y^{A-1})]$, hence
$q_{k_i}\bmod d\in G$.

The important point is that the random indices $k_1,\dots,k_K$ are determined
only by the good/block indicators, i.e. by the events ``exactly one prime is
selected in $\mathcal B_k$''. They are not determined by the residue classes of
the unique selected primes. More formally, condition on any possible value of
the index vector $(k_1,\dots,k_K)$. Since the blocks are disjoint, the selected
prime in each conditioned-good block is independent of the selected primes in
the other conditioned-good blocks, and the law inside each such block is the
law of the unique selected prime conditional on that block being good.

Equivalently, if $C_1,\ldots,C_K$ are arbitrary subsets of $G$, then on
conditioning on the event that the first $K$ good block indices are precisely
$k_1<\cdots<k_K$, independence of the disjoint prime blocks gives
\[
 \PP\bigl(q_{k_i}\bmod d\in C_i\ (1\le i\le K)\mid k_1,\ldots,k_K\bigr)
 =
 \prod_{i=1}^K
 \PP\bigl(q_{k_i}\bmod d\in C_i\mid \mathcal B_{k_i}\text{ is good}\bigr).
\]
The information that the intervening blocks are not good concerns only those
intervening blocks and therefore does not affect the conditional distribution
inside the good blocks $\mathcal B_{k_i}$.

By Lemma~\ref{lem:block-eq}, for each $i$ this conditional residue distribution
is $O(y^{-1/2})$-close in total variation to the uniform distribution on
$G=(\ZZ/d\ZZ)^\times$, uniformly in the realized index $k_i$. Thus, conditional
on the realized index vector, the residues
\[
g_i:=q_{k_i}\bmod d\qquad (1\le i\le K)
\]
are independent and each is $\eps_y$-close to uniform on $G$, where
\[
\eps_y:=O(y^{-1/2}).
\]
This conclusion is uniform over the conditioned index vector
$(k_1,\dots,k_K)$ by Remark~\ref{rem:block-uniformity}.

\medskip\noindent
\textbf{Step 6: Apply the subset-product lemma.}
The following estimate is first taken conditional on the event $W\ge K$ and on
an arbitrary realized index vector $k_1<\cdots<k_K$ for the first $K$ good
blocks. By the preceding step and Remark~\ref{rem:block-uniformity}, under this
conditioning the residues
\[
   g_i:=q_{k_i}\bmod d \qquad (1\le i\le K)
\]
are independent and each has total-variation distance $O(y^{-1/2})$ from the
uniform distribution on $G$. The bound obtained below is uniform in the realized
index vector, so it remains valid after averaging over the possible realized
vectors.

With Lemma~\ref{lem:subset-product} (and Remark~\ref{rem:near-uniform}) applied
to the $K$ residues $g_i:=q_{k_i}\bmod d\in G$,
\begin{align}
\PP\bigl(\nexists\,\emptyset\ne S\subseteq[K]:\;\textstyle\prod_{i\in S}g_i=1_G\bigr)
&\le \frac{N}{2^K-1} + K\eps_y\nonumber\\
&\le \frac{y}{2^{\kappa\log y - 1}} + O\bigl((\log y)y^{-1/2}\bigr)\nonumber\\
&= y^{1-\kappa\log 2}\cdot 2 + O\bigl(y^{-1/2}(\log y)\bigr)\nonumber\\
&= O(y^{-c_2}), \label{eq:subset-prob}
\end{align}
The estimate above is conditional on the event $W\ge K$ and on any fixed
realization of the first $K$ good block indices. Since the bound is uniform in
that realization, averaging over the possible index vectors preserves the same
bound.
where one may take
\[
c_2=\min\bigl(\kappa\log 2-1,\,\tfrac13\bigr)>0
\]
by \eqref{eq:gamma-choice}, since
$(\log y)y^{-1/2}=O(y^{-1/3})$.

\medskip\noindent
\textbf{Step 7: Putting it together.} Combining Steps 4 and 6,
\begin{equation}\label{eq:total-prob}
\PP(\text{no such product $e$ exists}) \;\le\; O(y^{-c_1})+O(y^{-c_2})\;=\;O(y^{-c}),
\qquad c:=\min(c_1,c_2)>0.
\end{equation}

When the construction succeeds, the resulting product
$e=\prod_{i\in S}q_{k_i}$ uses at most $K=\lceil\kappa\log y\rceil$ primes.
Since $\kappa=\rho\gamma/2<\gamma$ and $y$ is large, we have
$K\le \gamma\log y$. Each prime factor is at most $\exp(y^{A-1})$, so
\begin{equation}\label{eq:e-size}
\log e
\;\le\;
\sum_{i\in S}\log q_{k_i}
\;\le\;
K\cdot y^{A-1}
\;\le\;
\gamma(\log y)\cdot y^{A-1}
\;\le\;
y^A,
\end{equation}
since $(\log y)\le y$ for $y\ge 1$. Also $e>\exp y>y\ge d$ since $S\ne\emptyset$
and each prime in the product exceeds $\exp y$. The product is squarefree
because the primes $q_{k_i}$ in distinct good blocks are distinct (the blocks
are disjoint). Finally $e\equiv 1\pmod d$ by the subset-product condition.
Hence $e$ satisfies \eqref{eq:composite-successor}. \qed
\end{proof}

\section{Lower bound for $H(n)$}\label{sec:lower-H}

We now prove the missing matching lower bound:

\begin{theorem}\label{thm:lower-H}
For all but $o(x)$ integers $n\le x$,
\begin{equation}\label{eq:thm-lower-H}
H(n) \;\ge\; \bigl(1 - o(1)\bigr)\logstar x.
\end{equation}
\end{theorem}

\subsection{Tower scale for $H$}

Throughout this section we again write
\[
   L:=\logstar x.
\]

Let $A>10$ be the constant from Lemma~\ref{lem:composite-successor}, and choose
\begin{equation}\label{eq:B-choice}
B \;:=\; A+10.
\end{equation}
Set
\begin{equation}\label{eq:m0-mj-H}
m_0 \;:=\; \lfloor L^{1/2}\rfloor, \qquad y_j \;:=\; \exp\!\Bigl(\tfrac{T_{m_0+j-1}}{B}\Bigr) \quad (j\ge 1),
\end{equation}
and
\begin{equation}\label{eq:R-H}
R \;:=\; L - m_0 - 4.
\end{equation}
Note that $R = L - O(L^{1/2}) = (1-o(1))L$.

\begin{lemma}[Scale property for $H$]\label{lem:scale-H}
For all sufficiently large $L$ and all $1\le j\le R$,
\begin{equation}\label{eq:scale-H}
\exp(y_j^A) \;\le\; y_{j+1}.
\end{equation}
\end{lemma}

\begin{proof}
Set $m:=m_0+j-1$. Then
\[
y_j=\exp(T_m/B),\qquad y_{j+1}=\exp(T_{m+1}/B).
\]
The desired inequality is equivalent to
\[
y_j^A \;\le\; \log y_{j+1}=\frac{T_{m+1}}{B},
\]
or, since $T_{m+1}=\exp(T_m)$,
\[
\exp\!\Bigl(\frac{A T_m}{B}\Bigr)
\;\le\;
\frac{\exp(T_m)}{B}.
\]
Taking logarithms, this is
\[
\frac{A T_m}{B} \;\le\; T_m-\log B,
\]
equivalently
\[
T_m\Bigl(1-\frac{A}{B}\Bigr)\;\ge\; \log B.
\]
Because $B=A+10$, we have $1-A/B=10/B>0$. Also $m\ge m_0$ for every
$1\le j\le R$, and $m_0=\lfloor L^{1/2}\rfloor\to\infty$. Hence
$T_m\to\infty$ uniformly in $1\le j\le R$, so the displayed inequality holds
for all sufficiently large $L$. \qed
\end{proof}

\subsection{The greedy chain construction in the product model}

\begin{lemma}\label{lem:greedy-H}
Work in the independent prime model on primes $r$ up to $\exp(y_R^{A-1})$. Then with probability $1-o(1)$, there is a chain of integers
\begin{equation}\label{eq:H-chain}
1=d_1<d_2<\cdots<d_{R+1}, \qquad d_{j+1}\equiv 1\pmod{d_j},\qquad d_{j+1}\le y_{j+1},
\end{equation}
such that each $d_j$ ($j\ge 2$) is a squarefree product of selected primes lying in the disjoint windows $\bigl(\exp y_{j-1},\exp(y_{j-1}^{A-1})\bigr]$.
\end{lemma}

\begin{proof}
We construct the chain recursively and make deterministic choices. Set
$d_1:=1$. Suppose that $d_j$ has already been constructed, $d_j\le y_j$, and
$d_j$ is determined by the selections in the earlier windows $W_1,\dots,W_{j-1}$.

At stage $j$, use only primes in the fresh window
\begin{equation}\label{eq:windowj}
W_j \;:=\; \bigl(\exp y_j,\,\exp(y_j^{A-1})\bigr].
\end{equation}
Call a squarefree product $e$ \emph{admissible at stage $j$} if all its prime
factors are selected primes from $W_j$ and
\[
e\equiv 1\!\pmod{d_j},\qquad e>d_j,\qquad e\le \exp(y_j^A).
\]
If at least one admissible $e$ exists, define $d_{j+1}$ to be the least such
admissible integer; otherwise the construction fails at stage $j$.

At this point $d_j$ is measurable with respect to the selections in the earlier
windows $W_1,\dots,W_{j-1}$. The fresh window $W_j$ is disjoint from all earlier
windows, so conditional on the past and on the realized value of $d_j$, the
selection variables for primes in $W_j$ still have the independent Bernoulli
law required by Lemma~\ref{lem:composite-successor}. Since $d_j\le y_j$ and
Lemma~\ref{lem:composite-successor} is uniform for all $1\le d\le y_j$, it gives
\[
\PP(\text{stage $j$ fails}\mid\text{past})\;\ll\; y_j^{-c}.
\]
The deterministic choice of the least admissible $e$ ensures that $d_{j+1}$ is
measurable with respect to the selections in $W_1,\dots,W_j$.

On success, Lemma~\ref{lem:composite-successor} guarantees
\[
d_{j+1}\equiv 1\!\pmod{d_j},\qquad d_{j+1}>d_j,\qquad d_{j+1}\le\exp(y_j^A).
\]
By Lemma~\ref{lem:scale-H}, valid for every $1\le j\le R$,
$\exp(y_j^A)\le y_{j+1}$, so $d_{j+1}\le y_{j+1}$.

The windows $W_1,W_2,\dots,W_R$ are pairwise disjoint. Indeed,
Lemma~\ref{lem:scale-H} gives $y_{j+1}\ge\exp(y_j^A)>y_j^{A-1}$, hence
$\exp(y_j^{A-1})<\exp(y_{j+1})$, so the upper endpoint of $W_j$ lies below
the lower endpoint of $W_{j+1}$. Thus the primes used at stage $j$ are
independent of all previous choices.

Let $\mathcal F_j$ be the event that stage $j$ fails, and let
$\mathcal S_{j-1}$ be the event that all previous stages have succeeded.
The estimate above gives, uniformly in the realized past on $\mathcal S_{j-1}$,
\[
   \PP(\mathcal F_j\mid \text{past}) \ll y_j^{-c}.
\]
Therefore
\[
\PP(\mathcal F_j\cap \mathcal S_{j-1})
=
\EE\!\left[
   \mathbf 1_{\mathcal S_{j-1}}
   \PP(\mathcal F_j\mid \text{past})
\right]
\ll y_j^{-c}.
\]
Thus summing over $j$ is legitimate even though the modulus $d_j$ is random and
depends on the earlier stages.

Finally, by the union bound,
\begin{equation}\label{eq:greedy-H-bound}
\PP(\text{some stage fails})
\;\le\;
\sum_{j=1}^R O(y_j^{-c}).
\end{equation}
For all sufficiently large $L$, the sequence $y_j$ grows at least
quadratically: $y_{j+1}\ge y_j^2$. Hence
\[
\sum_{j=1}^R y_j^{-c}
\;\le\;
\sum_{j\ge 1} y_1^{-c\,2^{j-1}}
\;\le\;
2\,y_1^{-c}
\;=\;
o(1),
\]
because $y_1=\exp(T_{m_0}/B)\to\infty$. Therefore the construction succeeds
with probability $1-o(1)$. \qed
\end{proof}

\subsection{Proof of Theorem~\ref{thm:lower-H}}

The construction in Lemma~\ref{lem:greedy-H} only uses primes up to
\begin{equation}\label{eq:P-def-H}
P \;:=\; \exp\bigl(y_R^{A-1}\bigr)
\;=\; \exp\!\Bigl(\exp\!\Bigl(\tfrac{(A-1)T_{m_0+R-1}}{B}\Bigr)\Bigr).
\end{equation}
Now $m_0+R-1=L-5$, so
\[
   \log\log P=\frac{(A-1)T_{L-5}}{B}.
\]
Since $B=A+10$, we have $(A-1)/B<1$. Hence, for all sufficiently large $L$,
\begin{equation}\label{eq:P-bound1}
   \log\log P
   =
   \frac{(A-1)T_{L-5}}{B}
   <
   T_{L-5}.
\end{equation}
Exponentiating twice gives
\[
   \log P<T_{L-4}, \qquad P<T_{L-3}.
\]
By \eqref{eq:logx-vs-T}, $\log x>T_{L-2}\gg T_{L-3}$, so
\begin{equation}\label{eq:P-bound2}
P \;\le\; T_{L-3}\;=\;o(\log x).
\end{equation}

Set $M:=\prod_{p\le P}p$. By Chebyshev (Lemma~\ref{lem:Cheb}),
$\log M\le C_\theta P=o(\log x)$, hence $M=x^{o(1)}=o(x)$.

We also record explicitly that the constructed divisors lie below $x$. For
$1\le j\le R$,
\[
d_{j+1}\le y_{j+1}\le y_{R+1}
=
\exp\!\Bigl(\frac{T_{m_0+R}}{B}\Bigr)
=
\exp\!\Bigl(\frac{T_{L-4}}{B}\Bigr)
<
T_{L-3}
<
T_{L-2}
<
\log x
<
x
\]
for all sufficiently large $L$, using \eqref{eq:logx-vs-T}. Combined with
Remark~\ref{rem:CRT-products}, this shows that after CRT transfer each
constructed $d_j$ is a genuine divisor of the realized integer $n\le x$, not
merely a formal product in the independent model.

The event ``the greedy chain $\mathcal{C}_H$ of Lemma~\ref{lem:greedy-H} exists''
depends only on which primes $p\le P$ divide $n$. More explicitly, every window
$W_j$ used in the construction lies below $P$, every admissible divisor
$d_{j+1}$ is a squarefree product of selected primes from such a window, and the
rule choosing the least admissible product is deterministic once the indicators
$\bigl(\mathbf 1_{p\mid n}\bigr)_{p\le P}$ are known. No information about primes
above $P$, and no information about higher powers of primes, is used.

Although the event is phrased in terms of the existence of admissible
squarefree products, it is a finite event once the cutoff $P$ is fixed: each
window $W_j$ contains only primes $\le P$, and the admissible products are
subsets of this finite set of primes. The deterministic rule choosing the least
admissible product therefore makes the entire construction a well-defined
function of the finite vector
$\bigl(\mathbf 1_{p\mid n}\bigr)_{p\le P}$.

Hence by
Lemma~\ref{lem:CRT} its density among $n\le x$ equals
$\PP_{\mathrm{prod}}(\mathcal{C}_H)+O(M/x)=1-o(1)$.
Whenever $\mathcal{C}_H$ holds, the constructed divisors
\[
1=d_1<d_2<\cdots<d_{R+1}
\]
all divide $n$ and satisfy $d_{j+1}\equiv 1\pmod{d_j}$. Thus $H(n)\ge R+1$,
and in particular
\[
H(n)\;\ge\; R\;=\;(1-o(1))L.
\]
This proves \eqref{eq:thm-lower-H}. \qed

\section{Proof of Theorem~\ref{thm:main}}\label{sec:proof-of-main}

Combine Theorems~\ref{thm:upper-H},~\ref{thm:upper-h},
\ref{thm:lower-h}, and~\ref{thm:lower-H}. Thus, with $L:=\logstar x$, for all
but $o(x)$ integers $n\le x$,
\begin{align}
\bigl(\tfrac12-o(1)\bigr)L
\;\le\;
h(n)
&\;\le\;
\bigl(\tfrac12+o(1)\bigr)L,
\label{eq:h-L-bounds}\\
\bigl(1-o(1)\bigr)L
\;\le\;
H(n)
&\;\le\;
\bigl(1+o(1)\bigr)L.
\label{eq:H-L-bounds}
\end{align}
Here and below we have intersected the four density-one sets supplied by those
theorems; the union of their exceptional sets still has size $o(x)$.

It remains to replace $L=\logstar x$ by $\logstar n$. We first record a
uniform comparison on the interval $[x^{1/2},x]$. Since $\logstar x=L$, we
have
\[
T_{L-1}<x\le T_L
\]
for $L\ge 1$. Hence, for $n\in[x^{1/2},x]$,
\[
n\;\le\; x\;\le\; T_L,
\]
so $\logstar n\le L$. Also
\[
n \;\ge\; x^{1/2} \;>\; T_{L-1}^{1/2} \;=\; \exp(T_{L-2}/2).
\]
For all sufficiently large $L$, $\exp(T_{L-2}/2)>T_{L-2}$. Therefore $n>T_{L-2}$,
and so $\logstar n\ge L-1$. Consequently,
\begin{equation}\label{eq:logstar-comparison}
\logstar n=L+O(1)
\end{equation}
uniformly for $n\in[x^{1/2},x]$.

This comparison is uniform on $[x^{1/2},x]$. Hence any estimate of the form
$A(n)=cL+o(L)$ holding outside an exceptional set of size $o(x)$ may be
rewritten on the same density-one set, after discarding $n\le x^{1/2}$, as
\[
   A(n)=\bigl(c+o(1)\bigr)\logstar n,
\]
with the same constant $c$.

The integers $n\le x^{1/2}$ form an exceptional set of size $O(x^{1/2})=o(x)$.
Discarding them, \eqref{eq:logstar-comparison} gives
\[
L=(1+o(1))\logstar n
\]
uniformly on the remaining density-one set. Substituting into
\eqref{eq:h-L-bounds} and \eqref{eq:H-L-bounds} yields
\[
h(n)=\Bigl(\tfrac12+o(1)\Bigr)\logstar n,
\qquad
H(n)=\bigl(1+o(1)\bigr)\logstar n
\]
for all but $o(x)$ integers $n\le x$. Finally,
\[
\frac{H(n)}{h(n)}
\;=\;
\frac{\bigl(1+o(1)\bigr)\logstar n}
{\bigl(\tfrac12+o(1)\bigr)\logstar n}
\;=\;
2+o(1),
\]
because the lower bound for $h(n)$ ensures $h(n)\to\infty$ on the same
density-one set. This proves \eqref{eq:h-normal}--\eqref{eq:ratio}. \qed

\section{Concluding remarks}\label{sec:remarks}

\subsection{What is and is not proven} We have proven the
\emph{normal orders} of $h(n)$ and $H(n)$, i.e.\ asymptotic statements about a
density-one set of $n$. Theorem~\ref{thm:main} does \emph{not} assert that
$h(n)\le (\tfrac12+o(1))\logstar n$ pointwise, nor that $H(n)\ge (1-o(1))\logstar n$
pointwise; both pointwise statements are easily seen to fail (e.g.\ for $n$ a
product of a long prime chain $p_1\equiv 1\pmod{p_0}, p_2\equiv 1\pmod{p_1},\dots$).
Erd\H{o}s' original belief that $h(n)$ has a normal order $\logstar n$ is
\emph{at least partially} correct: we get $\Theta(\logstar n)$ with the precise
constant $1/2$.

\subsection{Sharpness of the second-order term} Our bounds carry an additive
slack of order $L^{1/2}$, where $L=\logstar x$, coming from the small-prime
cutoff $Y=\lfloor L^{1/2}\rfloor$. After replacing $L$ by $\logstar n$ on the
final density-one set, this is an additive
$O((\logstar n)^{1/2})$ slack. The truth is presumably
$h(n)=\tfrac12\logstar n+O(\log\logstar n)$ a.a., but proving this would require
a finer treatment of bad pairs at moderate scales. We have not attempted it
here.

\subsection{Convergence in probability}
The methods of Sections~\ref{sec:lower-h} and~\ref{sec:lower-H}, together with
the upper bounds, show convergence in probability:
\[
\frac{h(n)}{\logstar n}\to \frac12,
\qquad
\frac{H(n)}{\logstar n}\to 1
\]
with respect to natural density. We do not claim $L^p$-convergence here; that
would require additional uniform-integrability estimates beyond the arguments
of this paper.

\subsection{Why the constant is $1/2$ and not, say, $1$}
The factor of $2$ between the prime-chain and divisor-chain rates comes from
the fact that two prime steps are needed to pass through one full ``$U$-tower''
level, while one divisor step suffices. The bad-pair analysis is perfectly
parallel for $h$ and $H$, but the cost-per-step asymmetry --- which can be
traced to the
Brun--Titchmarsh-controlled scale $Q(p)=\exp\bigl(\exp\bigl(p/(\log p)^2\bigr)\bigr)$
versus the pure exponential scale $\exp\sqrt d$ --- forces the factor of $2$.

\subsection{Possible refinements}
\begin{enumerate}[label=(\alph*),leftmargin=*]
\item \emph{Sharper constants in second order.} One can probably push the
additive $O(L^{1/2})$ error in \eqref{eq:h-final-bound}, equivalently
$O((\logstar n)^{1/2})$ on the final density-one set, down to
$O(\log\logstar n)$ by being more careful with the bad-pair tail. We have not
explored this.
\item \emph{Distributional questions.} It is natural to ask whether, after
centering and normalizing,
$\bigl(h(n)-\tfrac12\logstar n\bigr)$ has a limiting distribution. This appears
to be a substantially harder question: the prime-chain construction in
Section~\ref{sec:lower-h} is essentially a renewal process, but the
Siegel--Walfisz error term is too crude to give a CLT directly.
\item \emph{Worst-case bounds.} The pointwise extremal problem is quite
different from the normal-order problem considered here. For instance, if
$n=\operatorname{lcm}(1,2,\ldots,k)$, then
\[
   1<2<3<\cdots<k
\]
is a valid divisor chain, since $j+1\equiv 1\pmod j$, and therefore
$H(n)\ge k$ while $\log n\sim k$ by the prime number theorem in the form
$\log\operatorname{lcm}(1,\ldots,k)\sim k$. Similarly, special integers built
from long prime chains can make $h(n)$ much larger than its normal order.
These pointwise phenomena are far from the density-one behavior proved here.
\end{enumerate}

\begin{thebibliography}{99}
\bibitem[Dav]{Davenport}
H. Davenport, \emph{Multiplicative Number Theory},
Graduate Texts in Mathematics \textbf{74}, 3rd edition revised by H.~L.~Montgomery,
Springer-Verlag, New York, 2000.

\bibitem[EGPS90]{EGPS1990}
P. Erd\H{o}s, A. Granville, C. Pomerance and C. Spiro,
\emph{On the normal behavior of the iterates of some arithmetic functions},
in \emph{Analytic Number Theory} (Allerton Park, IL, 1989),
Progress in Mathematics \textbf{85}, Birkh\"auser Boston, MA, 1990, pp.~165--204.

\bibitem[FKL10]{FKL2010}
K. Ford, S.~V. Konyagin and F. Luca,
\emph{Prime chains and Pratt trees},
Geometric and Functional Analysis \textbf{20} (2010), 1231--1258.

\bibitem[Erd]{ErdosProblems696}
The Erd\H{o}s Problems Database, problem \#696,
\url{https://erdosproblems.com/696} (consulted 2026-04-26).

\bibitem[HW]{HardyWright}
G.~H. Hardy and E.~M. Wright,
\emph{An Introduction to the Theory of Numbers},
6th edition, edited by D.~R.~Heath-Brown and J.~H.~Silverman,
Oxford University Press, Oxford, 2008.

\bibitem[IK]{IwaniecKowalski}
H. Iwaniec and E. Kowalski,
\emph{Analytic Number Theory},
American Mathematical Society Colloquium Publications \textbf{53},
American Mathematical Society, Providence, RI, 2004.

\bibitem[MV]{MontgomeryVaughan}
H.~L. Montgomery and R.~C. Vaughan,
\emph{The large sieve},
Mathematika \textbf{20} (1973), 119--134.

\bibitem[Mer]{Mertens-orig}
F. Mertens,
\emph{Ein Beitrag zur analytischen Zahlentheorie},
J. Reine Angew. Math. \textbf{78} (1874), 46--62.

\bibitem[SW]{SiegelWalfisz}
A. Walfisz,
\emph{Zur additiven Zahlentheorie. II},
Math. Z. \textbf{40} (1936), 592--607.

\end{thebibliography}

\end{document}
